\subsection{Singular image curves and the normalization conductor}

\begin{statement}{A.5.5}{Proved}{Conductor-cleared jets on a complete-intersection critical curve}
Let $Y$ be a smooth compact complex threefold with
$H^2(Y,\mathbb Q)=0$, let $S\subset Y$ be a reduced divisor,
and let $D\subset\operatorname{Sing}(S)$ be a reduced
irreducible compact curve which is a local complete intersection
in $Y$. Assume that the generic plane-curve transversal to $D$
in $S$ is an ordinary cusp, an ordinary tacnode, or an ordinary
triple point.
If the normalization $C\to D$ has positive genus, then $D$ is
a smooth elliptic curve. The transversal type is the indicated
cusp, tacnode, or triple point at every point of $D$;
in particular no degeneration of that type occurs along $D$.

In the cusp case the proof constructs, on $C$, the actual
meromorphic cubic on the radical of the Hessian and bounds its
pole divisor by twice the conductor. In the tacnode case the
corresponding fourth-order invariant has poles bounded by four
times the conductor. These intermediate jets need not extend
holomorphically before the conductor correction.
\end{statement}
\begin{proof}
Let $\nu:C\to D$ be the normalization and $j:C\to Y$
its composite with the inclusion,
$I=\mathcal I_{D/Y}$, and
\[
 N=\bigl(\nu^*(I/I^2)\bigr)^*,\qquad
 M=j^*\mathcal O_Y(S).
\]
These are a rank-two vector bundle and a line bundle on $C$.
Let $\Xi$ be the conductor divisor on $C$, so that
$\nu_*\mathcal O_C(-\Xi)$ is the conductor ideal in
$\mathcal O_D$, and put $c=\deg\Xi\geq0$.
Adjunction for the complete-intersection curve and finite
duality give exact line-bundle comparisons
\[
 \det N=\nu^*\omega_D\otimes j^*K_Y^{-1},\qquad
 \omega_C=\nu^*\omega_D\otimes\mathcal O_C(-\Xi).
                                                        \tag{A.J.1}
\]
For the second comparison one can work directly in the total
quotient ring: $\operatorname{Hom}_{\mathcal O_D}
(\nu_*\mathcal O_C,\mathcal O_D)$ consists precisely of
elements multiplying the normalization into $\mathcal O_D$,
which is its conductor. Since $D$ is Gorenstein, tensoring
this identity by its invertible dualizing sheaf is finite
duality. Both $M$ and $j^*K_Y$ have degree zero, by the
stated rational cohomology hypothesis. Consequently
\[
 \deg\det N=2g(C)-2+c.                              \tag{A.J.2}
\]
This is the normal bundle of the complete-intersection ideal,
with its specified normalization pullback. It is the bundle
needed in the following jet calculation.

First treat an ordinary cusp. Choose a local defining section
$f$ of $S$, using a local frame of $\mathcal O_Y(S)$.
Since $D\subset\operatorname{Sing}(S)$, its conormal class
maps to zero under
$I/I^2\to\Omega_Y^1|_D$. This homomorphism is injective:
its source is locally free on the reduced curve and the map
is injective at every generic point. Therefore $f\in I^2$.
The associated-graded identity for a regular sequence gives
$I^2/I^3=\operatorname{Sym}^2(I/I^2)$. The class of $f$
is thus a quadratic form
\[
 q\in H^0\bigl(C,\operatorname{Sym}^2N^*\otimes M\bigr).
\]
It has rank one generically. Let $R\subset N$ be its
saturated radical and $Q=N/R$. Both are line bundles.
The quadratic form annihilates $R$ everywhere by holomorphic
continuation, and gives a nonzero section of $M\otimes Q^{-2}$.
In particular
\[
 \deg Q\leq0.                                      \tag{A.J.3}
\]

On $D_{\rm reg}$, the third-order term on the radical is
intrinsic. Indeed, after writing the quadratic form as
$a\ell^2$, subtract a holomorphic lift of $a\ell^2$ from
$f$ and restrict its cubic class to $R^3$. Changing that lift
adds terms divisible by $\ell$, which vanish on $R^3$.
Changing normal generators has the same effect. Its nonzero
value characterizes the ordinary-cusp transversal among the
rank-one quadratic cases. Denote this cubic by $k$.

We now extend precisely this section across $C\setminus
\nu^{-1}(D_{\rm reg})$. Work at a point of $D$ and choose
two generators $X,Y$ of $I$. After a constant invertible linear
change of these generators, the normalized linear factor of
$q$ has the form $X+hY$ on every normalization branch over
that point. Such a choice exists: the saturated quotient
$N\to Q$ is surjective at each of the finitely many points,
so a constant vector can be chosen outside their finitely
many kernels. Thus $h$ is holomorphic on the normalization,
and
\[
 q=a(X+hY)^2,
\]
where $a=q_{XX}$ is the pullback of an actual coefficient
in $\mathcal O_D$. Choose a local conductor generator $d$
on the normalization. Both $d$ and $dh$ belong to
$\mathcal O_D$, since the conductor is an ideal in the
normalization contained in that ring. Choose lifts
$\widehat d,\widehat{dh},\widehat a$ to $\mathcal O_Y$,
and set
\[
 G=\widehat d^{\,2} f-
    \widehat a\bigl(\widehat d X+\widehat{dh}Y\bigr)^2.
                                                        \tag{A.J.4}
\]
The quadratic class of $G$ is zero after normalization.
It is already zero on $D$: the map from the locally free
sheaf $\operatorname{Sym}^2(I/I^2)$ to its normalization
pullback is injective because $D$ is reduced. Hence
$G\in I^3$.

Its class in $I^3/I^4=\operatorname{Sym}^3(I/I^2)$,
pulled back to $C$ and evaluated on $R^3$, is holomorphic.
Dividing that class by $d^2$ gives $k$ on the dense regular
locus. To check this equality, there $d$ is a unit and
(A.J.4) is exactly the subtraction of the square used in the
definition of the radical cubic. Different choices of lifts
add cubic terms containing its linear factor, and therefore
have zero restriction to $R^3$. Thus these local expressions
are the same meromorphic section wherever they overlap.
Since $d$ has divisor $\Xi$, the resulting bound is
\[
 0\ne k\in H^0\bigl(C,M\otimes R^{-3}
                               \otimes\mathcal O_C(2\Xi)\bigr).
                                                        \tag{A.J.5}
\]
This proves the claimed pole control without a splitting of
the formal neighbourhood of the singular curve.

Taking degrees in (A.J.5) gives $\deg R\leq2c/3$.
Together with (A.J.2)--(A.J.3), this yields
\[
 2g(C)-2+c=\deg R+\deg Q\leq\frac{2c}{3},
 \qquad 2g(C)-2+\frac c3\leq0.                     \tag{A.J.6}
\]
For $g(C)\geq1$, this forces $g(C)=1$ and $c=0$.
Conductor zero means $\mathcal O_D=\nu_*\mathcal O_C$;
a normal complex curve is smooth. Hence $D=C$ is elliptic.
Now (A.J.2)--(A.J.5) give $\deg R=\deg Q=0$.
The Hessian coefficient and radical cubic are nonzero
sections of degree-zero line bundles on this elliptic curve,
so neither has a zero. Locally the holomorphic Morse lemma
in the nondegenerate normal direction, with the curve
coordinate as parameter, writes $f=y^2+g(x,z)$.
The conditions just proved give $g=x^3u(x,z)$ with $u$ a
holomorphic unit. Replacing $x$ by a local cube root of $u$
times $x$ gives the actual local equation $y^2+x^3=0$.
This proves the asserted constant transversal type.

Next treat an ordinary tacnode. Retain $N,M,R,Q,q,a,h,d,P$
from the cusp calculation, with
$P=\widehat d X+\widehat{dh}Y$ and
$G=\widehat d^{\,2}f-\widehat aP^2\in I^3$.
The generic radical cubic is now zero, so its continuation
$[G]_3|_{R^3}$ is zero everywhere on the normalization.
Since $\ell=X+hY$ is monic in $X$, polynomial division gives
an actual holomorphic quadratic $C_2$ on the normalization
such that
\[
 [G]_3=\ell C_2.
\]
Each coefficient of $dC_2$ is in the conductor and hence in
$\mathcal O_D$. Lift those coefficients and let $H$ be the
resulting homogeneous quadratic in $X,Y$. Set
\[
 K=\widehat d^{\,2}G-PH.
\]
Its cubic class pulls back to
$d^2\ell C_2-d\ell\cdot dC_2=0$.
The same locally free graded-piece injection used above gives
$K\in I^4$. There is an exact identity in the local ring of
the ambient smooth threefold:
\[
 4\widehat a\widehat d^{\,6}f
 -(2\widehat a\widehat d^{\,2}P+H)^2
 =4\widehat a\widehat d^{\,2}K-H^2.               \tag{A.J.7}
\]
On the dense smooth locus, complete the square in the
nondegenerate normal direction. If the relevant terms there
are $a\ell^2+b\ell r^2+er^4$, the fourth-order invariant is
$\Delta=ae-b^2/4$. A substitution
$\ell\mapsto\ell+ur^2$ leaves this expression unchanged;
terms divisible by $\ell^2$ and higher changes in the radical
coordinate do not alter it. Thus it is the intrinsic section
of $M^2\otimes Q^{-2}\otimes R^{-4}$ on that dense locus.
It is nonzero precisely for the ordinary-tacnode transversal
among the rank-one Hessian cases with zero radical cubic.

Divide (A.J.7) by $4\widehat d^{\,6}$ and take the
fourth-order class on $R^4$. Since
$[H]_2=dC_2$, this expresses exactly that same invariant as
\[
 \Delta=
 \frac{a[K]_4|_{R^4}-(C_2|_{R^2})^2/4}{d^4}.
                                                        \tag{A.J.8}
\]
Its numerator is holomorphic on the normalization.
Different choices agree on the dense regular locus and hence
give one meromorphic section. Therefore
\[
 0\ne\Delta\in H^0\bigl(C,M^2\otimes Q^{-2}
                    \otimes R^{-4}\otimes\mathcal O_C(4\Xi)\bigr).
                                                        \tag{A.J.9}
\]
Write $q_0=\deg Q\leq0$. Using
$\deg R=2g(C)-2+c-q_0$, the degree of the line bundle
in (A.J.9) is
\[
 -2q_0-4\bigl(2g(C)-2+c-q_0\bigr)+4c
       =-8g(C)+8+2q_0.
\]
Its nonzero section and $g(C)\geq1$ force
$g(C)=1$ and $q_0=0$.

The section $q$ of $M\otimes Q^{-2}$ is now nowhere zero.
In the local presentation above, $a=q_{XX}$ is therefore a
unit on every normalization branch. Since it belongs to
$\mathcal O_D$, it is a unit in that local ring as well.
The coefficient $h=q_{XY}/a$ consequently belongs to
$\mathcal O_D$; here $q_{XY}$ denotes the coefficient in
the symmetric matrix, so the cross term is $2q_{XY}XY$.
We can thus repeat the entire construction with $d=1$.
Polynomial division by the monic $X+hY$ then takes place
over $\mathcal O_D$ itself; the zero radical cubic descends
by the same injectivity argument. Formula (A.J.8) now has
no poles. The resulting nonzero section of
$M^2\otimes Q^{-2}\otimes R^{-4}$ has degree $-4c$,
which forces $c=0$.

Thus again $D=C$ is smooth elliptic. Both the nonzero
quadratic coefficient and $\Delta$ are sections of
degree-zero line bundles and are nowhere zero. The parametric
Morse lemma first gives $f=y^2+g(x,z)$, and the vanished
radical cubic and nonzero fourth invariant give
$g=x^4u(x,z)$ with $u$ a unit. Taking a local fourth root
produces the actual equation $y^2+x^4=0$ at every point of
$D$. This proves the tacnode assertion.

For an ordinary triple transversal, the quadratic class
vanishes generically, hence vanishes in the locally free
sheaf $I^2/I^3$. Thus $f\in I^3$, and its initial cubic
is an honest global section
$b\in H^0(C,\operatorname{Sym}^3N^*\otimes M)$.
The discriminant of a binary cubic is homogeneous of degree
four in its coefficients and transforms by the sixth power
of the determinant of a change of variables. Its value is
therefore a section
\[
 0\ne\operatorname{disc}(b)\in
 H^0\bigl(C,M^4\otimes(\det N)^{-6}\bigr).
                                                        \tag{A.J.10}
\]
It is nonzero because the three generic tangent lines are
distinct. Equations (A.J.2) and (A.J.10) give
$2g(C)-2+c\leq0$, again forcing $g(C)=1$ and $c=0$
under the genus assumption. The discriminant has degree
zero and is nowhere zero. The three simple tangent roots
extend holomorphically on each small curve disc. The implicit
function theorem gives three smooth local branches of $S$
with those pairwise distinct normal tangents. Choose the
first two branches as the coordinate axes; the third is
$y=xu(x,z)$ with $u$ a unit, and replacing $y$ by $y/u(x,z)$
straightens it to $y=x$. This gives the local equation
$xy(x-y)=0$ up to a unit at every point of $D$.
\end{proof}


The pole allowance in (A.J.5) has a concrete local use. Put
$h=z^2-x^3$, $D=(h,w)=0$ in $\mathbb C^3$, and
\[
 f=xw^2+2zwh+x^2h^2+h^3.
\]
On the normalization $x=s^2,z=s^3$ of $D$, its quadratic
class is $s^2(w+sh)^2$. Completing the square on $D_{\rm reg}$
gives
\[
 f=x\left(w+\frac zx h\right)^2+
                  \left(1-\frac1x\right)h^3.
\]
The radical cubic has a pole of order two at $s=0$.
The conductor generator is $d=x=s^2$, and the exact
conductor-cleared identity is
\[
 x^2f-x(xw+zh)^2=x(x-1)h^3\in I_D^3.
\]
Thus the corrected jet comparison applies even in this local
model where holomorphic extension of the uncorrected cubic
would fail. The equalities are checked exactly in
\texttt{code/E\_singular\_curve.py}.

\subsection{A double conductor with smooth normalization}

The local calculation below is related to the conductor
derivations and characteristic-zero pole calculation in
Reid~\cite[\S3.9, \S4.4--4.8]{Reid1994Nonnormal}.
We give the canonical-module comparison directly, without
using an anticanonical polarization or a classification of
the surface.

\begin{statement}{A.5.7}{Proved}{A pure double conductor forces a smooth cusp image}
Let $S$ be a reduced Gorenstein complex surface and
$\pi:\overline S\to S$ its finite normalization.
Suppose that, over a neighbourhood of a point $p\in S$,
$\overline S$ is smooth and its conductor ideal is exactly
$\mathcal O_{\overline S}(-2C)$ for a smooth reduced
divisor $C$; its local branches over $p$ are disjoint.
Suppose that the generic transversal of $S$ along the
reduced image $D=\pi(C)$ is an ordinary cusp.
Then $D$ is smooth at $p$, $C\to D$ is an isomorphism
near the preimage of $p$, and the germ of $S$ at $p$ is
analytically a product of an ordinary cusp and a smooth disc.
No complete-intersection hypothesis on $D$ is required.
\end{statement}
\begin{proof}
Write $R=\mathcal O_{S,p}$ and let $\overline R$ be its
finite semilocal normalization. It is a product of regular
two-dimensional local analytic rings. Write $\mathfrak c$
for the conductor, and set
\[
 A=R/\mathfrak c,\qquad
 \overline A=\overline R/\mathfrak c,\qquad
 Q=\overline R/R=\overline A/A.
\]
We first give the exact duality comparison needed below.
The $R$-module $\overline R$ is maximal Cohen--Macaulay.
Since $R$ is Gorenstein,
$\operatorname{Ext}_R^j(\overline R,R)=0$ for $j>0$.
Its dual $\operatorname{Hom}_R(\overline R,R)$ is
precisely the conductor, viewed as a fractional ideal:
a homomorphism is multiplication in the total quotient ring,
and multiplication takes $\overline R$ into $R$ exactly
for conductor elements. Applying $\operatorname{Hom}_R(-,R)$
to $0\to R\to\overline R\to Q\to0$ therefore gives
\[
 \operatorname{Ext}_R^1(Q,R)=A.                    \tag{A.J.11}
\]
The depth lemma shows that $Q$ is a Cohen--Macaulay module
of dimension one. Gorenstein duality on such modules is the
involutive functor $M\mapsto\operatorname{Ext}_R^1(M,R)$.
Consequently $A$ is Cohen--Macaulay and
\[
 Q\simeq\operatorname{Ext}_R^1(A,R)=\omega_A,
 \qquad \operatorname{End}_A(Q)=A.                \tag{A.J.12}
\]
The final equality is the canonical multiplication map:
duality identifies its endomorphism ring with
$\operatorname{End}_A(A)=A$.

The ordinary-cusp hypothesis says that $A$ is generically
reduced. Its nilradical is therefore supported at finitely
many points; a one-dimensional Cohen--Macaulay ring has
no nonzero submodule with that support. Hence $A$ is reduced.
Choose local coordinates $(s_i,z_i)$ on the normalization
branches with $C=(z_i=0)$. Put
\[
 B=\prod_i\mathbb C\{s_i\},\qquad
 K=\prod_i\operatorname{Frac}\bigl(\mathbb C\{s_i\}\bigr).
\]
The specified conductor gives the actual ring comparison
\[
 \overline A=B[\epsilon]/(\epsilon^2),             \tag{A.J.13}
\]
component by component. The projection
$\overline A\to B$ is injective on $A$, since its kernel
is nilpotent and $A$ is reduced. It identifies $B$ with
the finite normalization of the curve ring $A$.

Under (A.J.13) the inclusion of $A$ takes the form
\[
 a\longmapsto a+\delta(a)\epsilon,
\]
where $\delta:A\to B$ is a complex-linear analytic
derivation. It extends to the total quotient ring, and on
each normalization branch has the form
$\delta(a)=u_i\,da/ds_i$ for a meromorphic germ $u_i$.
Let $u=(u_i)$ and write derivatives componentwise.
The map
\[
 \Phi:B[\epsilon]/(\epsilon^2)\longrightarrow K,
 \qquad b+c\epsilon\longmapsto c-u b'
                                                        \tag{A.J.14}
\]
is $A$-linear for this specified action: multiplication by
$a+u a'\epsilon$ changes its value to
$ac+u a'b-u(ab)'=a(c-u b')$. Its kernel at every generic
point is exactly the graph of the extended derivation on
$K$. It therefore induces a generic isomorphism from $Q$
to its image in $K$. The module $Q\simeq\omega_A$ is
torsion-free on the reduced curve, so this induced map is
injective.

Differentiation of convergent power-series germs onto
convergent power-series germs is surjective, by termwise
integration. Hence the image in (A.J.14) is exactly
\[
 I=B+uB\subset K.                                \tag{A.J.15}
\]
It is an invertible fractional ideal on the semilocal
normalization $B$. Every $A$-linear endomorphism of this
rank-one module is multiplication by an element of $K$;
preserving $I$ is equivalent, branch by branch, to belonging
to $B$. Thus $\operatorname{End}_A(I)=B$ as subrings of
$K$. Combining this with the actual $A$-linear isomorphism
$Q\xrightarrow{\Phi}I$ and (A.J.12) gives $A=B$.
Because $A$ is local, there is a single branch over $p$,
and its ring is regular. This proves that $D$ is smooth
and that $C\to D$ is an isomorphism there.

Finally $A=B=\mathbb C\{s\}$ makes $u=\delta(s)$
holomorphic. The analytic coordinate change
$s'=s+u(s)z$ has invertible differential at $z=0$ and
identifies the image of $A$ modulo $z^2$ with the functions
of $s'$ alone. Since $R$ is the inverse image of $A$ under
$\overline R\to\overline A$, it follows that
\[
 R=\mathbb C\{s'\}+z^2\mathbb C\{s',z\}
   =\mathbb C\{s',z^2,z^3\}
   \simeq\mathbb C\{s',x,y\}/(y^2-x^3).
                                                        \tag{A.J.16}
\]
This is the asserted product germ.
\end{proof}

\subsection{Vertical exceptional curves on a ruled resolution}

\begin{statement}{A.5.8}{Proved}{Rational singularities for a vertically resolved ruled normalization}
Let $\pi:Z\to T$ be a resolution of a compact connected
normal complex surface, where $Z$ is an irrational ruled
surface with ruling $p:Z\to B$. Suppose that every
exceptional curve of $\pi$ is vertical for $p$.
Then there is a unique holomorphic map
$\widetilde p:T\to B$ satisfying
$p=\widetilde p\circ\pi$, and
\[
 R\pi_*\mathcal O_Z=\mathcal O_T,
 \qquad R^1\widetilde p_*\mathcal O_T=0.           \tag{A.J.17}
\]
In particular all singularities of $T$ are rational. For
$g=g(B)$, its complete structure-sheaf cohomology dimensions
are
\[
 (h^0(T,\mathcal O_T),h^1(T,\mathcal O_T),
                         h^2(T,\mathcal O_T))=(1,g,0).
                                                        \tag{A.J.18}
\]
\end{statement}
\begin{proof}
The exceptional fibre of the proper birational map $\pi$
over a point is connected, because $T$ is normal and
$\pi_*\mathcal O_Z=\mathcal O_T$. Its curve components
are vertical by assumption, so their intersection
connectivity makes $p$ constant on that whole fibre.
It is also constant on the one-point fibres. We verify
holomorphic descent, rather than only descending the
underlying continuous map.

Fix $t\in T$ and let $U\subset B$ be a coordinate disc
containing $p(\pi^{-1}(t))$. Properness makes
$\pi(Z\setminus p^{-1}(U))$ closed and disjoint from
$t$. There is therefore an open neighbourhood $W$ of $t$
with $\pi^{-1}(W)\subset p^{-1}(U)$. The coordinate
function of $p$ on $\pi^{-1}(W)$ is a section of
$\pi_*\mathcal O_Z(W)=\mathcal O_T(W)$, so descends
to a holomorphic function on $W$. Its values lie in the
coordinate image of $U$; the inverse coordinate gives a
holomorphic map $W\to U$. These local maps agree after
pullback to the surjective map $\pi$, hence agree on
overlaps. They construct the required unique
$\widetilde p$.

Contracting vertical ruling $(-1)$-curves factors $p$ as
\[
 Z\xrightarrow{b} Z_0\xrightarrow{p_0}B,
\]
where $Z_0\to B$ is a $\mathbb P^1$-bundle and $b$
is a finite sequence of point blowups. Every such blowup
has direct image $\mathcal O$ and no higher direct image
of $\mathcal O$: the successive layers on its
infinitesimal exceptional curves are
$\mathcal O_{\mathbb P^1}(k)$ for $k\geq0$, whose
degree-one cohomology is zero, and formal functions gives
the assertion. The bundle $p_0$ has the same direct-image
property by the cohomology of $\mathbb P^1$ and local
triviality. Composition consequently gives
\[
 p_*\mathcal O_Z=\mathcal O_B,
 \qquad R^j p_*\mathcal O_Z=0\quad(j>0).          \tag{A.J.19}
\]

The sheaf $\mathcal T=R^1\pi_*\mathcal O_Z$ has
finite support at the singular points of $T$. The
low-degree Leray sequence for $p=\widetilde p\circ\pi$
contains the actual arrows
\[
 0\longrightarrow R^1\widetilde p_*\mathcal O_T
 \longrightarrow R^1p_*\mathcal O_Z
 \longrightarrow\widetilde p_*\mathcal T
 \longrightarrow R^2\widetilde p_*\mathcal O_T.    \tag{A.J.20}
\]
The last term is zero: the nonconstant proper map
$\widetilde p$ has fibres of complex dimension one.
The middle term is zero by (A.J.19). Thus
$R^1\widetilde p_*\mathcal O_T=0$ and
$\widetilde p_*\mathcal T=0$. Pushforward is faithful
on a sheaf with finite support, since its stalk sections
at finitely many points cannot disappear under a proper
map. Hence $\mathcal T=0$. Higher $R^j\pi_*$ vanish
for $j>1$ by the dimension of its fibres. This proves
(A.J.17) and the defining rational-singularity property.

Finally (A.J.17) identifies every
$H^q(T,\mathcal O_T)$ with $H^q(Z,\mathcal O_Z)$.
Equation (A.J.19) identifies the latter with
$H^q(B,\mathcal O_B)$. A compact connected smooth curve
has dimensions $(1,g)$ in degrees zero and one and no
higher cohomology, proving (A.J.18).
\end{proof}

\subsection{The two derivative directions on a ruled normalization}

\begin{statement}{A.5.10}{Proved}{Eliminating vertical conductor on a smooth ruled cusp normalization}
Let $Y$ be a smooth compact complex threefold with
$H^2(Y,\mathbb Q)=0$, let $S\subset Y$ be a reduced
irreducible divisor, and suppose its normalization
$h:Z\to S\subset Y$ is a smooth irrational
ruled surface $p:Z\to B$.
Suppose its conductor divisor is $D_Z=2C+V$, where
$C$ is a ruling section and $V$ is vertical, and
the generic transversal along $h(C)$ is an ordinary cusp.
Then $B$ is elliptic, $V=0$, $C^2=0$, and $Z$ is
already minimal.
The image $h(C)$ is smooth elliptic and $S$ is locally
an ordinary-cusp product along it. No local
complete-intersection hypothesis on the image curve is
required. The conductor-image scheme along the vertical
components may initially be nonreduced.
\end{statement}
\begin{proof}
Let $g=g(B)\geq1$, let $F$ be the full ruling fibre
class, and write $V=\sum_i c_i\Gamma_i$ for its
irreducible vertical components. Put
$v=V\cdot C$, $n=C^2$ and $k=2g-2$.
Conductor duality and Cartier adjunction give the exact
line-bundle comparison
\[
 K_Z\otimes\mathcal O_Z(2C+V)
       =h^*(K_Y\otimes\mathcal O_Y(S)).
                                                        \tag{A.J.21}
\]
Its right side has rational first Chern class zero.
Intersecting with $C$ and using adjunction on the section
gives
\[
 n=2-2g-v=-k-v.                                  \tag{A.J.22}
\]
Let $D_p$ be the divisorial common zero of
$dp:T_Z\to p^*T_B$. Its remaining zero scheme is
zero-dimensional, say $\Xi$, so the image is
$I_\Xi\otimes p^*T_B(-D_p)$.
The kernel $T_p$ is an invertible sheaf, by the depth
lemma. Taking determinants gives the actual comparison
\[
 T_p=K_Z^{-1}\otimes p^*K_B\otimes\mathcal O_Z(D_p).
\]
The section makes $dp$ surjective at every point of $C$.
Consequently $D_p$ and $\Xi$ are disjoint from $C$,
and $T_p|_C=N_{C/Z}$.

Put $\mathcal F=h^*T_Y$. Poincare duality on the compact
oriented six-manifold $Y$ also gives
$H^4(Y,\mathbb Q)=0$. Therefore
\[
 c_1(\mathcal F)=0,\qquad c_2(\mathcal F)=0
        \quad\hbox{in rational cohomology}.       \tag{A.J.23}
\]
The differential is generically injective. Let
$L\subset\mathcal F$ be the saturation of the image
of $T_p\xrightarrow{dh}\mathcal F$.
It is invertible: a saturated rank-one subsheaf of a
locally free sheaf on a smooth surface is reflexive.
The map $T_p\to L$ has no vertical divisorial zero.
Indeed $h$ is finite, so no ruling fibre is contracted;
in characteristic zero its derivative along that fibre
is nonzero generically.

At the generic point of $C$, an ordinary cusp has
normalization differential equivalent to that of
$(u,s)\mapsto(u^2,u^3,s)$.
The ruling tangent either lies in its kernel along $C$
or does not. The vanishing order in the first case is
exactly one, and in the second case it is zero.
Write $\epsilon=1$ or $\epsilon=0$ for these respective
alternatives. There are no other horizontal conductor
components, and the normalization is an immersion away
from the conductor. Hence the actual line comparison is
\[
 L=T_p\otimes\mathcal O_Z(\epsilon C).
                                                        \tag{A.J.24}
\]
This checks both positions of the derivative kernel
relative to the ruling tangent.

The quotient $E=\mathcal F/L$ is torsion-free of rank
two. The differential induces a generically nonzero map
\[
 T_Z/T_p=I_\Xi\otimes p^*T_B(-D_p)
                    \longrightarrow E\longrightarrow E^{**}.
\]
It extends uniquely to a map from $p^*T_B(-D_p)$ into
$E^{**}$, since both are locally free and the omitted
set has codimension two; locally this is Hartogs
extension of its coefficient functions.
Let $M$ be its saturated image line in the rank-two
bundle $E^{**}$. At the generic point of $C$, the
two-by-two differential minors have common vanishing
order one. Removing the order $\epsilon$ already
assigned to $L$ leaves order $1-\epsilon$ for this map.
At the generic point of $\Gamma_i$, take a plane transversal
to its image and a normalization parameter $u$ for the
corresponding plane-curve branch. If $m_i$ is its
multiplicity, its derivative modulo the nonzero tangent
along $\Gamma_i$ vanishes to order $r_i=m_i-1$.
These are all codimension-one zeros. Put
$R=\sum_i r_i\Gamma_i$ and $r=R\cdot C$. We obtain
\[
 M=p^*T_B(-D_p)\otimes
           \mathcal O_Z\bigl((1-\epsilon)C+R\bigr).
                                                        \tag{A.J.25}
\]

The plane-curve conductor gives the inequalities
\[
 2R\leq V,\qquad 2r\leq v.                      \tag{A.J.26}
\]
To spell out their local length content, on a single
plane branch of multiplicity $m_i$ the classes
$u,u^2,\ldots,u^{m_i-1}$ are linearly independent in
the normalization quotient. Its delta invariant is thus
at least $m_i-1$. The branch is Gorenstein. Finite
duality identifies the length of its ring modulo its
conductor with the length of its normalization quotient;
the intrinsic conductor exponent is consequently twice
that delta invariant. Other branches add their
nonnegative intersection multiplicities to this exponent.
For the last comparison, writing the plane equation
as $\prod_j f_j$ multiplies the residue denominator on
branch $i$ by $\prod_{j\ne i}f_j$, whose valuation is
exactly the sum of those intersection multiplicities.
This proves (A.J.26) for every branch, including cases
with nonreduced conductor-image scheme.

Let $\ell_1$ be the length of $E^{**}/E$. Saturation
of $M$ gives a zero-dimensional subscheme $\Xi_M$ and a
line bundle $N$ with
\[
 0\longrightarrow M\longrightarrow E^{**}
       \longrightarrow I_{\Xi_M}\otimes N\longrightarrow0.
\]
Write $\ell_2=\operatorname{length}\Xi_M$.
First Chern classes give $c_1(N)=-c_1(L)-c_1(M)$.
The second Chern classes in these typed quotient
sequences and (A.J.23) give the intersection identity
\[
 L^2+L\cdot M+M^2=\ell_1+\ell_2\geq0.            \tag{A.J.27}
\]
Indeed $c_2(E)=L^2$,
$c_2(E^{**})=L^2-\ell_1$, and
$c_2(E^{**})=-L\cdot M-M^2+\ell_2$.

Set $W=V+D_p$ and $B_v=R-D_p$, both vertical
divisors, the first effective. The exact comparisons
above and (A.J.21) give the numerical classes
\[
 L\equiv(2+\epsilon)C+W+kF,\qquad
 M\equiv(1-\epsilon)C+B_v-kF.
\]
The intersection form on vertical divisors is negative
semidefinite. One can check this directly by contracting
the ruling blowups: every vertical class is a multiple
of the pullback fibre class plus exceptional classes,
on which the form is negative definite. Hence
\[
 q_v=W^2+W\cdot B_v+B_v^2
      =(W+B_v/2)^2+3B_v^2/4\leq0.
\]
Its vanishing implies $W^2=B_v^2=0$; these classes
are then numerical multiples of the fibre class.

For $\epsilon=1$, using $D_p\cdot C=0$ in the
left side of (A.J.27) gives
\[
 \ell_1+\ell_2=-6k-3v+3r+q_v
                  \leq-6k-3v/2.
\]
Nonnegativity, with $k\geq0$ and $v\geq0$,
forces $g=1$, $v=0$, $r=0$ and $q_v=0$.

For $\epsilon=0$, the same intersection calculation reads
\[
 0\leq\ell_1+\ell_2
       =-6k-2v+4r+q_v\leq-6k.                   \tag{A.J.28}
\]
It forces $g=1$, $\ell_1=\ell_2=0$, $2r=v$,
and $q_v=0$.
Thus $E=E^{**}$ is locally free, and
$L\to\mathcal F$ is a subbundle inclusion.
Along $C$, the composite
$T_C\to\mathcal F|_C\to E|_C$ vanishes generically,
because $dh$ has rank one there and its nonzero ruling
tangent image is $L|_C$. It therefore vanishes
everywhere. The exact subbundle sequence gives an
actual nonzero holomorphic line-bundle map
\[
 T_C\longrightarrow L|_C=T_p|_C=N_{C/Z}.
                                                        \tag{A.J.29}
\]
It is nonzero since the finite map $C\to h(C)$ has
nonzero derivative generically. Its Hom line has degree
$n-(2-2g)=-v$ by (A.J.22). A nonzero section on a
compact curve has nonnegative degree, so $v=0$ in this
alternative too.

In both alternatives, $v=r=0$ and $q_v=0$.
Thus $W$ is numerically a fibre multiple and
$W\cdot C=V\cdot C+D_p\cdot C=0$, so it is
numerically zero. An effective nonzero divisor on the
projective ruled surface has positive intersection
with an ample class. It follows that $W=0$ and hence
$V=D_p=0$. Now $B_v=R$ is effective, has square zero
and has zero intersection with $C$, so the same
argument gives $R=0$. Equation (A.J.22) gives $C^2=0$.
Finally $K_Z+2C$ is numerically trivial. A ruling
$(-1)$-curve $E_0$ would satisfy
$0=(K_Z+2C)\cdot E_0=-1+2C\cdot E_0$, impossible
because $C\cdot E_0$ is an integer. There are therefore
no ruling blowups and $Z$ is minimal.

Statement A.5.7 now applies everywhere along
the conductor: the normalization is smooth and its
conductor is exactly $2C$. It makes the image smooth
and the surface locally the cusp product. The image is
elliptic because $C\to h(C)$ is an isomorphism.
\end{proof}


\subsection{An unconditional conductor genus bound}

\begin{statement}{A.5.14}{Proved}{Arithmetic genus two and the embedded delta-one comparison}
Let $f:Y\to C_0$ be a proper connected-fibre map from a
smooth compact complex threefold with $H^2(Y,\Q)=0$ to
a smooth curve. Suppose its smooth general fibre has
trivial canonical bundle. Let $mS$ be a full fibre with
irreducible reduced nonnormal support, and let $D$ be
a reduced irreducible component of its conductor image.
Then
\[
                         p_a(D)\le2.                \tag{A.J.30}
\]
If the normalization of $D$ has positive genus, then
$D$ is a local complete intersection in $Y$; every
singular germ of $D$ is an ordinary node or an ordinary
cusp, and the sum of its delta invariants is at most
one. In particular, when the generic transverse type
along such a positive-genus $D$ is $A_2$, $A_3$, or an
ordinary triple point, $D$ is smooth elliptic and the
transverse analytic type is constant along all of $D$.
These conclusions do not require a nonconstant ambient
Picard pullback.
\end{statement}
\begin{proof}
The evaluation, adjunction and fibre-equation comparisons
in (A.CD.1) use no Picard assumption and apply with
rational second-cohomology vanishing. They give a torsion
dualizing line $\omega_S$. The divisor Riemann--Roch
calculation gives $\chi(\OO_S)=0$. Thus, for
$e=h^0(S,\omega_S)\in\{0,1\}$, the same Serre-duality
calculation as in (A.CD.2) gives
$h^1(\OO_S)=1+e$; if $e=1$, then
$\omega_S\simeq\OO_S$.

Write $\eta:T\to S$ for normalization,
$\mathcal Q=\eta_*\OO_T/\OO_S$, and $\Gamma$ for
the full conductor image scheme. Equality of the
constant global functions on $S$ and $T$ gives the
actual injection $H^0(\mathcal Q)\hookrightarrow
H^1(\OO_S)$. Conductor duality and curve Serre duality
therefore give
\[
 h^1(\Gamma,\omega_S|_\Gamma)
       =h^0(S,\mathcal Q)\le1+e.                  \tag{A.J.31}
\]
The restriction quotient
$\omega_S|_\Gamma\twoheadrightarrow\omega_S|_D$
has a kernel supported in dimension at most one, so
its $H^1$ map is surjective. For $e=1$ it follows that
$p_a(D)=h^1(\OO_D)\le2$. For $e=0$, put
$L=\omega_S|_D$. This torsion line has degree zero,
and (A.J.31) gives $h^1(D,L)\le1$. Curve
Riemann--Roch gives
\[
 p_a(D)=1+h^1(D,L)-h^0(D,L)\le2,
\]
proving (A.J.30) in both cases.

Let $\nu:\widehat D\to D$ be normalization. The
normalization sequence gives the exact numerical
comparison
\[
 p_a(D)=g(\widehat D)+
        \sum_{x\in D}\dim_{\C}
          (\nu_*\OO_{\widehat D}/\OO_D)_x.       \tag{A.J.32}
\]
If $g(\widehat D)>0$, every local delta is consequently
at most one and their sum is at most one. We give the
analytic local rings at a point with delta one.
Let $A=\OO_{D,x}$ and let
$B=\bigoplus_{i=1}^r\C\{t_i\}$ be its finite
normalization. Here $\dim_{\C}(B/A)=1$. The residue map
of the local ring $A$ has the same value on every
normalization branch. The resulting quotient of
$\C^r$ by its diagonal has dimension $r-1$, so
$r\le2$.

For $r=2$, the subring
\[
 A_0=\{(b_1,b_2)\in\C\{t_1\}\oplus\C\{t_2\}:
                                  b_1(0)=b_2(0)\}
\]
contains $A$ and has colength one in $B$. Hence
$A=A_0$. The explicit homomorphism
$\C\{u,v\}/(uv)\to A_0$ sends
$u$ to $(t_1,0)$, $v$ to $(0,t_2)$, and a constant
to the corresponding equal pair. Expanding a
convergent series shows that it is an isomorphism.
This is the ordinary node germ.

For $r=1$, an element of $A$ with valuation one would
be a local coordinate on the normalization disc by
the holomorphic inverse function theorem. Convergent
composition in the analytic algebra $A$ would then
give $\C\{t\}\subset A$, contrary to its colength
one. Therefore every nonconstant element of $A$ has
valuation at least two. We have
$A\subset\C+t^2\C\{t\}$; both rings have
colength one in $\C\{t\}$, so equality holds.
The map
\[
 \C\{u,v\}/(v^2-u^3)\longrightarrow A,
                 \qquad (u,v)\longmapsto(t^2,t^3)
\]
is an isomorphism: every series with no linear term
is uniquely the sum of an even series in $t^2$ and
$t^3$ times a series in $t^2$. This is the ordinary
cusp germ. A delta-zero germ is its smooth
normalization disc.

We finally check the embedding into the specified
threefold. These analytic rings have embedding
dimension at most two. In local coordinates of the
smooth threefold, the tangent space of $D$ therefore
has codimension at least one. There is an element
$z$ of its ideal with $dz(x)\ne0$. The inverse
function theorem makes $z$ an ambient coordinate,
and the ideal of $D$ is the inverse image of a
height-one ideal in the regular analytic ring
$\C\{u,v\}=\OO_{Y,x}/(z)$. This ideal is principal:
a pure reduced curve in a smooth surface is a Cartier
divisor, or equivalently its reduced height-one ideal
is a product of the prime factors in this factorial
ring. Thus $I_{D,x}=(z,h)$, a regular sequence.
This proves the asserted lci embedding, including
at the node and cusp. Statement A.5.5 now applies
to the three specified generic transverse types
and proves the final smooth elliptic and constant-type
conclusion.
\end{proof}

\subsection{A local spherical-cycle obstruction}

\begin{statement}{A.5.15}{Proved}{Hirzebruch germs cannot occur in a torus degeneration}
Let $f:Y\to C_0$ be a proper holomorphic map from a
smooth complex threefold to a smooth curve, with smooth
nearby fibres complex two-tori. There is no point of
a fibre at which, in holomorphic ambient coordinates
$(x,y,z)$ and a coordinate on the base, the map has
the equation
\[
 f(x,y,z)=u(x,y,z)\bigl(x^a-y^bz^c\bigr)^m,
 \quad a\ge2,\quad b,c,m\ge1,                                    \tag{A.J.33}
\]
where $u$ is a holomorphic unit. No Betti-number
assumption on $Y$ is needed. This includes the pinch
point, every germ $x^2-y^{2k}z^{2r+1}$ with
$k\ge1$, $r\ge0$, and the reducible germ
$x^2-y^2z^2$.
\end{statement}
\begin{proof}
Put $h=x^a-y^bz^c$. We first compare the actual map
with its reduced equation. This polynomial is weighted
homogeneous of degree $a(b+c)$ for weights
$(b+c,a,a)$. On a small simply connected coordinate
neighbourhood choose a holomorphic logarithm of $u$.
Define the holomorphic coordinate change
\[
 \Phi(x,y,z)=\left(
 x u^{1/(ma)},\;
 y u^{1/(m(b+c))},\;
 z u^{1/(m(b+c))}\right).
                                                        \tag{A.J.34}
\]
Its derivative at the origin is an invertible diagonal
matrix. The holomorphic inverse function theorem and
weighted homogeneity give
$h\circ\Phi=u^{1/m}h$ and therefore
$f=(h\circ\Phi)^m$ exactly. For a sufficiently small
nonzero regular value $t$, a chosen Milnor fibre of
$h$ at one of the $m$ roots of $t$ is consequently a
compact oriented submanifold with boundary of the
smooth global fibre $Y_t$. Denote this Milnor fibre
by $M$.

We first give a direct spherical cycle, without an
irreducibility assumption. Put $d=\gcd(b,c)$,
$b'=b/d$ and $c'=c/d$, and consider the smooth affine
curve
\[
 B=\{(x,v)\in\C^2:x^a-v^d=1\}.
\]
The affine surface $H=\{h=1\}$ maps holomorphically
to $B$ by $v=y^{b'}z^{c'}$. The curve $B$ is
connected: projection to the $v$-line has $a$ sheets,
and a small loop about any root of $v^d+1$ cyclically
permutes them. It is smooth because its two partial
derivatives cannot vanish on $x^a-v^d=1$.
The function $v$ has $a$ distinct simple zeros
$P_j=(\zeta_a^j,0)$ on $B$.

Choose two of these zeros, $P$ and $Q$, and an
embedded arc $\gamma:[0,1]\to B$ joining them and
avoiding the other zeros. Choose smooth complex
functions $A(t),B_0(t)$ on its interior such that
\[
 A(t)^{b'}B_0(t)^{c'}=v(\gamma(t)),\qquad
 |A(t)|=|v(\gamma(t))|^{1/(2b')},\quad
 |B_0(t)|=|v(\gamma(t))|^{1/(2c')}.
\]
They exist by choosing a continuous argument along
the arc. The continuous map
\[
 (t,\theta)\longmapsto
 \bigl(x(\gamma(t)),A(t)e^{\mathrm i c'\theta},
                 B_0(t)e^{-\mathrm i b'\theta}\bigr)
                                                        \tag{A.J.33a}
\]
from $[0,1]\times S^1$ to $H$ collapses each boundary
circle to the corresponding point $P$ or $Q$.
It therefore defines an oriented spherical cycle
$\sigma_\gamma:S^2\to H$. Coprimality of $b'$ and
$c'$ makes each interior circle embedded. The whole
map is injective away from the collapsed ends,
because the base arc is embedded.

Let $\gamma'$ be a nearby arc with the same endpoints
and disjoint interior, obtained by an isotopy in a
thin disc neighbourhood of $\gamma$. The isotopy
keeps all other zeros of $v$ outside that disc.
Choose the arguments in (A.J.33a) continuously during
the isotopy; both coordinates tend uniformly to zero
at the ends. This gives a continuous homotopy of
maps from $S^2$, so the resulting spherical cycles
are homologous. Their supports meet only at $P$ and
$Q$. Near either endpoint, $(y,z)$ are actual
holomorphic coordinates on $H$, since $a x^{a-1}$
is nonzero there. Arrange the two arcs to be radial
with different arguments in the local $v$-coordinate.
On a sufficiently small sphere in these coordinates,
the links of the two cycles are disjoint parallel
curves of slope $(c',-b')$ on a Clifford torus.
Their oriented linking number is $-b'c'$.
Indeed the torus framing of a $(p,q)$ curve has
linking number $pq$ with its parallel copy: push the
two curves to the opposite sides of the Heegaard
splitting of $S^3$ into its two solid tori and count
meridian discs, obtaining $p q$ signed intersections.
Here $p=c'$ and $q=-b'$.

The local intersection of the two coned cycles is
that linking number. This equality follows directly
by replacing the two cones inside the small ball by
transverse relative two-chains with the same boundary
links; their intersection is the defining linking
number in its boundary three-sphere. At $Q$ both
cycle orientations reverse relative to the outgoing
radial coordinate, so the two reversals leave the
local intersection sign unchanged. Consequently
\[
 [\sigma_\gamma]^2
   =[\sigma_\gamma]\cdot[\sigma_{\gamma'}]
   =-2b'c'\ne0.                                  \tag{A.J.33b}
\]
Both cycles and their homotopy have compact support.
The weighted homothety
$(x,y,z)\mapsto(s^{b+c}x,s^a y,s^a z)$ maps
$H$ biholomorphically onto $\{h=s^{a(b+c)}\}$.
For sufficiently small positive $s$, it puts those
compact supports into any prescribed Milnor
neighbourhood of the origin. Thus (A.J.33b) is a
spherical class of nonzero square in an actual
local Milnor fibre $M$.

For an independent check when $\gcd(a,b,c)=1$,
we compute the full homotopy type and the rank of
the intersection form from two primary results.
The join theorem
\cite[Theorem L1 and Example 2.4(a)]{Nemethi1991Join}
applies to arbitrary analytic germs, including
non-isolated ones. Set $d=\gcd(b,c)$. The Milnor
fibre of $y^bz^c$ is a disjoint union of $d$ annuli:
its equations in a Milnor polydisc are
$y^{b/d}z^{c/d}=\zeta\epsilon^{1/d}$, and each
component is parametrized by
$y=A w^{c/d}$, $z=B w^{-b/d}$. Coprimality makes
this parametrization injective, and the two radius
bounds give a nonempty annulus in $w$.
The Milnor fibre of $x^a$ consists of $a$ points.
The join theorem thus identifies $M$, by a homotopy
equivalence, with the union of $a$ cones on a
disjoint union of $d$ circles, all glued along their
common base. Contracting one cone and then choosing
arcs between the $d$ base components gives
\[
 M\simeq
 \bigvee^{(a-1)d}S^2\ \vee\
 \bigvee^{(a-1)(d-1)}S^1.                          \tag{A.J.35}
\]
In particular, every class in $H_2(M,\Z)$ is a sum
of images of spherical fundamental classes.

For the same irreducible germ, the boundary calculation
of \cite[Theorem 4.1 and Lemma 4.6]{MPW2005Hirzebruch}
gives a Seifert three-manifold with base genus
\[
 g_\partial=(a-1)(d-1),\qquad
 e_0=\frac{a}{(b/d)(c/d)}>0.                     \tag{A.J.36}
\]
Here $e_0$ is its rational Seifert Euler number in
the orientation convention of that paper. Its
nonzero value gives
$b_1(\partial M)=2g_\partial$. To make this last
deduction explicit, the rational abelianized Seifert
presentation has a regular-fibre generator $v$,
exceptional-fibre generators $q_i$, the relations
$\alpha_iq_i+\beta_iv=0$ and
$\sum_iq_i+ev=0$, and $2g_\partial$ free base
generators. Eliminating the $q_i$ gives $e_0v=0$.
Thus all fibre generators vanish over $\Q$, leaving
exactly the $2g_\partial$ base generators.

Let $I_M:H_2(M,\Q)\to H_2(M,\partial M;\Q)$ be
the intersection morphism. Poincare--Lefschetz duality
and the homology sequence of the pair give
\[
 \operatorname{rank} I_M
   =b_2(M)+b_1(M)-b_1(\partial M)=a-1>0.          \tag{A.J.37}
\]
Indeed $H_3(M)=0$ by (A.J.35), so
$H_3(M,\partial M)\simeq H^1(M)$ injects into
$H_2(\partial M)$. Its cokernel is the kernel of
$I_M$, which has dimension
$b_1(\partial M)-b_1(M)$. This proves the displayed
rank formula. For $a=2$, $b=1$, $c=\ell$, the
more precise boundary is $L(2\ell,1)$ by
\cite[Theorem 4.2(2)]{MPW2005Hirzebruch}; the same
pair sequence over $\Z$ makes the rank-one
intersection form have absolute determinant
$2\ell$. The special case $\ell=1$ is the ordinary
quadratic singularity and has determinant two.

Finally let $j:M\hookrightarrow Y_t$ be the actual
embedding obtained from (A.J.34). A complex two-torus
has contractible universal cover $\C^2$, so
$\pi_2(Y_t)=0$. In particular the spherical class
constructed in (A.J.33a) maps to zero in
$H_2(Y_t,\Z)$. Intersection numbers of compact
two-cycles in the interior of $M$ are preserved by
this oriented codimension-zero embedding: the same
small transverse perturbations compute them in
both manifolds. Its square in $Y_t$ is therefore
both zero and $-2b'c'$, contradicting (A.J.33b).
In the irreducible case, (A.J.35)--(A.J.37) provide
an additional contradiction on the whole intersection
space.
This proves the exclusion for the actual map,
including every multiplicity in (A.J.33).
\end{proof}

The homotopy and boundary theorems just used are
established results; our application is the displayed
embedding and intersection comparison. The explicit
cycle also treats the reducible monomials. The exact
Seifert arithmetic, homology dimensions and ranks
are checked in \path{code/E_Hirzebruch.py}. No
reduction of an arbitrary nonordinary surface germ
to (A.J.33) is asserted.

\begin{statement}{A.5.17}{Proved}{Isolated singularities are excluded and suspension germs are transverse products}
Let $f:Y\to C_0$ have the hypotheses of A.5.15.
No local equation $f=u h^m$, with $u$ a holomorphic
unit and $m\ge1$, can have a reduced equation $h$
with an isolated critical point at the origin.
Suppose its local equation at a point of a fibre is
\[
             f=(x^a-g(y,z))^m,
             \qquad a\ge2,\quad m\ge1,             \tag{A.J.38}
\]
where $g\ne0$ and $g(0,0)=0$. Then the reduced
plane germ $\{g=0\}_{\rm red}$ is smooth. If
$b$ is its multiplicity in $\operatorname{div}(g)$,
there are local holomorphic coordinates in which
$f=(x^a-y^b)^m$.
Consequently, at every point where the full fibre
is a uniform multiple of its reduced support and
the reduced equation has a nonzero Hessian, its
critical surface germ is a constant
transverse $A_{b-1}$ product. A nonordinary change
of transverse type at a point with nonzero Hessian
is excluded.
\end{statement}
\begin{proof}
We first prove the isolated assertion, including its
comparison with the actual global fibre. Replace $h$
by $H=u^{1/m}h$, taking a holomorphic root on a small
coordinate neighbourhood. Then $f=H^m$ exactly.
This replacement preserves an isolated critical point
as a germ. Indeed every local irreducible component
of the critical locus of $H$ has constant $H$-value,
necessarily zero when it contains the origin. On
$H=0$ one has $dH=u^{1/m}dh$. Thus its critical
locus has the same zero-dimensional germ as that
of $h$.

Choose a small closed ball $K\subset\C^3$ such that
$H^{-1}(0)$ is smooth along its boundary intersection
and transverse to $\partial K$.
The standard isolated-singularity Milnor-ball
construction supplies this choice; see
\cite[Lemma 6.1 and Appendix B]{Milnor1968}.
The restriction $H|_{\partial K}$ is a submersion
at its zero set. Compactness therefore supplies
$\eta>0$ for which it is a submersion wherever
$|H|<2\eta$ on the boundary. For the linear
deformation
\[
 H_\lambda(w)=H(w)-\lambda_1w_1
                         -\lambda_2w_2-\lambda_3w_3,
                                                        \tag{A.J.38a}
\]
this boundary submersion persists for all parameters
in a sufficiently small polydisc $P$ and all values
of modulus less than $\eta$. The gradient map of
$H$ has an isolated zero and consequently is a finite
holomorphic map germ of positive local degree. Choose
a small regular value $\lambda$ of that map. Its
nonempty finite inverse image consists of critical
points of $H_\lambda$ with invertible Hessian. All
these points and their critical values tend to the
origin as $\lambda\to0$.

Remove the critical-value discriminant from
$P\times\Delta_\eta$. The relative critical locus
over this domain misses the boundary, so its
projection is proper and its image is a proper
complex analytic subset. Its complement is path
connected. The projection of the corresponding
smooth level sets in $K$ is a proper submersion
with boundary, and hence an Ehresmann bundle.
At any one of the nondegenerate critical points,
the holomorphic Morse lemma gives
$H_\lambda-H_\lambda(p)=X^2+Y^2+Z^2$.
A local Milnor fibre at a sufficiently small regular
value embeds in the corresponding level set in $K$.
The bundle transports it, by an oriented embedding,
into $K\cap H^{-1}(\epsilon_0)$ for a nonzero small
$\epsilon_0$, which is a submanifold of the actual
global torus $Y_{\epsilon_0^m}$.

For a positive real local Milnor value, the real
locus $X,Y,Z\in\mathbb R$ is a two-sphere. It is
Lagrangian for the standard K\"ahler form; multiplication
by $\mathrm i$ identifies its tangent bundle with
its normal bundle with the opposite oriented Euler
class. Its self-intersection is therefore
$-\chi(S^2)=-2$. Equivalently this is the case
$a=2$, $b=c=1$ of the explicit cycle in A.5.15,
after the linear change $Y^2+Z^2=(Y+\mathrm iZ)(Y-\mathrm iZ)$.
The transported sphere has the same square in the
actual torus, whose $\pi_2$ is zero. This is the
required contradiction. No suspension representation
of the original isolated singularity has been used.

We prove the comparison of Milnor fibres needed for
the argument. Factor the analytic plane function as
\[
                   g=u\prod_{i=1}^r h_i^{b_i},
\]
where the $h_i$ are distinct reduced irreducible
germs and $u$ is a unit. Deform each $h_i$
holomorphically to $h_{i,s}$ in a finite-dimensional
parameter polydisc $P$, and retain its exponent:
\[
 g_s=u\prod_i h_{i,s}^{b_i},\qquad
 H_s=x^a-g_s.                                    \tag{A.J.39}
\]
We use only sufficiently small parameters.

Choose a small closed ball $K\subset\C^2_{y,z}$
such that each $h_i=0$ meets $\partial K$
transversely in its smooth part, and their boundary
links are disjoint. Choose its representative so
that every critical value of $g$ on $K$ is zero.
This is possible because the critical locus of a
holomorphic function has finitely many local
irreducible components, the function is constant
on each of them, and each component of the germ
contains the origin. The boundary smoothness,
transversality and disjointness persist for
$h_{i,s}$ by compactness and the implicit function
theorem. Near any of those links, the other factors
are units, so there are parameter-dependent local
coordinates $(v,w)$ with $g_s=v^{b_i}$, obtained
by taking a root of that unit. The boundary
$\partial K$ is transverse to the $w$ direction.
Consequently the restriction of $g_s$ to
$\partial K$ is a submersion at every point where
$0<|g_s|<\eta_0$, for a uniform sufficiently small
$\eta_0>0$. Away from these link neighbourhoods
$|g_s|$ has a uniform positive lower bound on the
boundary. This proves the asserted uniform
transversality, including at parameters for which
the germs inside $K$ change their singularities.

Take a small disc $D_x$ of radius $\rho$ with
$\rho^a<\eta_0/2$, then shrink $P$ so that every
critical value of $g_s$ on $K$ has modulus less than
$\rho^a/4$. This shrinking follows by compactness:
a sequence of critical points with parameters
tending to zero would otherwise limit to a
nonzero critical value of $g$ in the chosen small
representative. The latter has no such value.
Fix $0<\eta<\rho^a/4$ and put
$W=\overline D_x\times K$, a compact manifold
with corners contained in the original coordinate
chart in $Y$.

Every $H_s^{-1}(\epsilon)$ with
$0<|\epsilon|<\eta$ is transverse to this boundary.
On the face $\partial D_x\times K$, the equation
makes $|g_s|>3\rho^a/4$, so $dg_s\ne0$; its plane
directions give the submersion. On the face
$\overline D_x\times\partial K$, the $x$
derivative gives it when $x\ne0$. At $x=0$ we
have $g_s=-\epsilon\ne0$, and the preceding
boundary calculation gives it. At the corner
$|x|=\rho$, that same calculation for $g_s$
supplies both real target directions. The corner
stratum is therefore also transverse. We use
Ehresmann's theorem with corners: local lifts of
base tangent vectors can be chosen tangent to
each boundary face by these submersion statements,
then patched by a partition of unity and
integrated on the compact fibres. Thus no change
of the chosen domain $W$ is required.

Inside $P\times\{0<|\epsilon|<\eta\}$ remove the
critical-value discriminant $\mathcal D$ of the
maps $H_s$ on $W$. This is a proper complex
analytic subset. The relative critical locus has
no point on the boundary at such a value, so its
projection is proper over every compact subset
of this parameter domain, and the proper mapping
theorem applies. For fixed $s$ its critical
values are discrete, since a holomorphic function
is constant on each connected critical component.
The complement of a proper complex analytic
subset of a connected complex manifold is path
connected. The transversality just proved and
the implicit function theorem give a proper
smooth submersion, including its boundary,
\[
 \{(s,p):p\in W,\ (s,H_s(p))\in
                    (P\times\Delta_\eta^*)\setminus\mathcal D\}
 \longrightarrow
       (P\times\Delta_\eta^*)\setminus\mathcal D,
 \quad (s,p)\longmapsto(s,H_s(p)).                 \tag{A.J.40}
\]
Ehresmann's theorem identifies all its compact
fibres by orientation-preserving diffeomorphisms
of manifolds with corners.
In particular a sufficiently small local Milnor
fibre of an interior singularity of some $H_s$,
at a value near zero, embeds into a smooth fibre
in (A.J.40), hence into
$W\cap H_0^{-1}(\epsilon_0)$ for a sufficiently
small nonzero $\epsilon_0$. This last space is
an actual submanifold of the global torus
$Y_{\epsilon_0^m}$ by (A.J.38); no identification
of it with a standard Milnor fibre is needed.
Choose the value to avoid every other
critical value, and choose its local Milnor ball
in the interior of $W$. This is the precise
comparison we use. Arbitrary deformations which
change the exponents along the boundary links
are not used in (A.J.40).

Suppose first that $r\ge2$. Choose two distinct
factors $h_i,h_j$. Their map
$(h_i,h_j):(\C^2,0)\to(\C^2,0)$ is finite,
since their common zero germ is zero dimensional.
Its positive local degree is their intersection
multiplicity. A sufficiently small regular value
$(s_i,s_j)$ consequently has nonempty finite
inverse image. At each inverse-image point,
$h_i-s_i$ and $h_j-s_j$ are transverse coordinates.
Choose the remaining small shifts to avoid those
finitely many points. The deformation (A.J.39)
then has an interior point at which
\[
            H_s=x^a-U v^{b_i}w^{b_j},
\]
with $U$ a unit independent of $x$. Replacing
$v$ by $vU^{1/b_i}$ makes this exactly
$x^a-v^{b_i}w^{b_j}$ by an invertible holomorphic
coordinate change. A.5.15 constructs a spherical
class of nonzero square in its local Milnor fibre.
The embedding (A.J.40) puts that class in the
actual smooth torus fibre of $f=H_0^m$.
Preservation of intersection
numbers and $\pi_2$-vanishing give the same
contradiction as in A.5.15. Thus $r=1$.

Suppose its reduced factor $h=h_1$ is singular.
Write $g=u h^b$ and absorb a local $b$th root of
$u$ into $h$. This leaves its reduced zero germ
unchanged. Since $h$ is reduced, its gradient
has an isolated zero at the origin: a critical
curve would have constant value zero and would
be a repeated component of $h$. The finite
gradient map has a positive local degree.
Choose a small regular value
$(\lambda_1,\lambda_2)$ of this map. Then
\[
        h_s=h-\lambda_1y-\lambda_2z-\lambda_0
\]
has nondegenerate critical points near zero;
regularity of the gradient value says exactly
that their Hessians are invertible. Choose
$\lambda_0$ so that their critical values are
nonzero. At one such point $p$ put
$v=h_s(p)\ne0$. The Hessian of $h_s^b$ is
$b v^{b-1}$ times the Hessian of $h_s$, and
is invertible. The holomorphic Morse lemma on
the two plane variables identifies the germ of
$H_s-H_s(p)$ with
\[
                         x^a-Y^2-Z^2.
\]
A linear change of $(Y,Z)$ gives the Hirzebruch
germ $x^a-YZ$, already excluded by A.5.15.
Its Milnor fibre embeds by (A.J.40), producing
the same contradiction. Hence $h$ is smooth.
Taking $y=h u^{1/b}$ as a coordinate proves the
asserted transverse product.

For the Hessian consequence, take a local $m$th
root of the actual map, absorbing its unit.
Its reduced equation has the same nonzero
Hessian rank as every other reduced equation
at the critical point. The holomorphic splitting
lemma splits off a nondegenerate square and
gives $f=(x^2-g(y,z))^m$ exactly. If $g=0$,
the equation would not be reduced, so $g\ne0$.
At a critical point $dg(0)=0$. The result just
proved gives $g=y^b$ with $b\ge2$ in suitable
holomorphic coordinates, which is the constant
transverse $A_{b-1}$ product.
\end{proof}

\subsection{The branch divisor of a horizontal tacnodal conductor}

\begin{statement}{A.5.18}{Proved}{Every horizontal tacnode image is elliptic}
Let $f:Y\to C_0$ be a proper holomorphic map from a smooth
compact complex threefold with $b_2(Y)=0$, with connected
fibres and smooth general fibre a complex two-torus.
Let $mS$ be a full fibre with irreducible reduced support,
$\eta:T\to S$ its normalization, and $\pi:Z\to T$ a
resolution with an irrational ruling $p:Z\to B$.
Suppose every exceptional curve of $\pi$ is vertical and
the conductor divisor on $Z$ is $2C+V$, where $C$ is a
section of $p$ and $V$ is vertical. Let $D=\eta\pi(C)$.
Suppose the generic plane-curve transversal to $D$ in $S$ is an
ordinary tacnode.

Then $D$ is a smooth elliptic curve, $B=C$ is elliptic,
and $C\to D$ is an unramified double cover. The transversal
type along $D$ is everywhere the same ordinary tacnode.
The normalization $T$ is a smooth minimal ruled surface
over this elliptic curve. Its conductor is twice a
section, and $e(S)=0$. In particular its minimal
resolution is the identity.
In particular a rational image of
the horizontal tacnodal conductor cannot occur. The proof
allows vertical conductor and rational singularities of
$T$; it does not assume that the normalization map is
immersive at the special points under examination.
\end{statement}
\begin{proof}
If the normalization of $D$ has positive genus, A.5.14
and A.5.5 make $D$ smooth elliptic and its transversal
type constant. There are two generic normalization
branches of a tacnode. A vertical component of $Z$
is rational, so it cannot dominate this elliptic curve.
Thus both branches lie on $C$, and $C\to D$ is an
unramified double cover: the constant local equation
$x^2-y^4=0$ has two separately smooth normalization
branches on every sufficiently small disc in $D$.
Riemann--Hurwitz gives $g(C)=g(B)=1$.
It remains to exclude the case $D^\nu=\Sigma\simeq\PP^1$.

The finite map $C\to\Sigma$ has degree at least two,
because $C\simeq B$ has positive genus. Since a generic
tacnode has exactly two normalization branches, this
degree is exactly two. Consequently no vertical
conductor component can also dominate $D$: its generic
points would supply an additional normalization branch.
Write $g=g(C)\ge1$ and let $\mathcal B\subset\Sigma$
be the reduced branch divisor of this double cover.
Riemann--Hurwitz gives
\[
                   \deg\mathcal B=2g+2.             \tag{A.T.1}
\]

We first bound the number of branches of the singular
image curve itself. A.5.8 supplies the descended ruling
$\widetilde p:T\to B$. Identify $C$ with $B$ using the
section, let $\phi:B\to D$ be its actual map to the
image curve, and put
\[
                    \psi=\phi\widetilde p:T\to D.
\]
Every vertical conductor component has constant $\psi$
value. Its image under $\eta$ differs from $D$, as just
proved, and meets $D$ in finitely many points. Thus
there is a finite set $A\subset D$ such that, for each
$s\in S$, either $\psi(\eta^{-1}(s))$ is a singleton
or all its values belong to $A$. Include in $A$ the
constant values on the vertical conductor components
and the values of $\phi$ over their finitely many
intersections with $D$. Fibres over the remaining
points of $D$ have their preimages on $C$, where all
the $\psi$ values equal that same point of $D$.

Choose an embedded finite tree $J\subset D$ containing
$A$, in a triangulation of the compact analytic curve.
It can be constructed by joining these finitely many
points by arcs, subdividing their intersections and
retaining a spanning tree. The quotient
$q:D\to D/J$ is a homotopy equivalence: a contractible
CW subcomplex is a cofibration, and the homotopy
extension property extends its contraction. Since
$q\psi$ is constant on every fibre of the finite
proper quotient $\eta$, there is a continuous map
\[
 a:S\to D/J,\qquad a\eta=q\psi,\qquad a|_D=q.     \tag{A.T.2}
\]
The last equality holds first on the dense regular
locus of $D$ and then everywhere by continuity.

Use the explicit triangulated neighbourhood retraction
$r:N\to S$ of (A.Pi.3). For a sufficiently small disc
$\Delta$ around the fibre value, properness gives
$U=f^{-1}(\Delta)\subset N$. The map $(ar)|_U$ is
surjective on fundamental groups because its restriction
to $D$ is the homotopy equivalence $q$. Perturbing loops
off the real-codimension-two analytic support gives
$\pi_1(U\setminus S)\twoheadrightarrow\pi_1(U)$.
The torus-bundle exact sequence over $\Delta^*$ therefore
gives the actual quotient maps
\[
 \Z^4\rtimes\Z\twoheadrightarrow\pi_1(U)
       \twoheadrightarrow\pi_1(D/J)\simeq\pi_1(D).
                                                        \tag{A.T.3}
\]
The first group is metabelian, so the last is metabelian.

For each singular point $x$ of $D$, let $r_x$ be the
number of its normalization branches and put
$\beta=\sum_x(r_x-1)$. The finite normalization map
identifies exactly those finite sets of points on
$\PP^1$; since it is a compact Hausdorff quotient,
this is also the actual underlying topological quotient.
Adjoining arcs between the points to be identified
gives a CW presentation homotopy equivalent to
$S^2\vee\bigvee^\beta S^1$. Hence
$\pi_1(D)$ is the free group $F_\beta$.
For $\beta\ge2$ this group surjects onto $S_4$ by
sending two free generators to $(1234)$ and $(12)$;
the commutator $A_4$ of that quotient is nonabelian.
Equation (A.T.3) therefore forces $\beta\le1$.
By A.5.14 the total delta invariant $\delta=p_a(D)$
is at most two. If $s$ denotes the number of points
of $\Sigma$ over $\operatorname{Sing}D$, then
\[
 s=\#\operatorname{Sing}D+\beta\le\delta+\beta\le3.
                                                        \tag{A.T.4}
\]

First suppose that $D$ is a local complete intersection.
We retain the exact bundles and jets of A.5.5.
Let $\nu:\Sigma\to D$ be normalization, $I=\mathcal I_D$,
\[
 N=(\nu^*(I/I^2))^*,\qquad
 M=\nu^*(\mathcal O_Y(S)|_D).
\]
Thus $M$ is a degree-zero line bundle on $\Sigma$.
Let $\Xi$ be the curve-normalization
conductor divisor, $c=\deg\Xi=2\delta$, and let
$R\subset N$ and $Q=N/R$ be the saturated radical and
quotient of the generic rank-one Hessian. Write
$q_0=\deg Q$. The adjunction and conductor comparisons
(A.J.1) give
\[
 \deg N=-2+c,\quad \deg R=-2+c-q_0,\quad
 0\ne a\in H^0(\Sigma,M\otimes Q^{-2}),\quad q_0\le0.
                                                        \tag{A.T.5}
\]
Here $a$ denotes the intrinsic quadratic coefficient,
whose zero divisor has degree $-2q_0$.
The actual quartic invariant $\Delta=ae-b^2/4$ of
A.5.5 is meromorphic in
$L^2$, where $L=M\otimes Q^{-1}\otimes R^{-2}$.
Its pole bound (A.J.9) gives a nonzero holomorphic
section $\sigma$ of the actual square line bundle
$L^2(4\Xi)=(L(2\Xi))^2$. Its zero divisor has degree
\[
               d_\Delta=8+2q_0.                    \tag{A.T.6}
\]

The double cover $C\to\Sigma$ is obtained by adjoining
$\sqrt{-\Delta}$ to the function field, with the line
trivializations specified by $L$. Indeed over the
generic point a tacnode has normal expansion
$a\ell^2+b\ell r^2+er^4$, and the two normalization
branches are distinguished by the two roots of
$a(\ell/r^2)^2+b(\ell/r^2)+e$. Their discriminant
is $-4\Delta$. Changing a local frame of $L$ multiplies
the representing function of $\Delta$ by a square,
so this describes the same degree-two extension on
every overlap. At a point with local parameter $z$,
write that representing function as $z^n u(z)$,
$u(0)\ne0$. A root of the unit and removal of the
even power of $z$ show that its normalization is
ramified exactly when $n$ is odd. Adding $4\Xi$ changes
each order by an even number. Thus $\mathcal B$ is
exactly the odd-multiplicity support of
$\operatorname{div}(\sigma)$, and
\[
                  2g+2\le 8+2q_0.                 \tag{A.T.7}
\]

At a point of $D_{\rm reg}$ where $a\ne0$, A.5.17
applies to the actual local $m$th root of the full-fibre
map. Its nonzero Hessian makes the local equation a
constant transverse product, necessarily of type
$A_3$ because this is the generic type along $D$.
The conductor cover is therefore unramified there.
Every branch point outside $\nu^{-1}(\operatorname{Sing}D)$
is consequently a zero of $a$. Equations (A.T.4)--(A.T.5)
give the second bound
\[
                    2g+2\le -2q_0+s.              \tag{A.T.8}
\]
Adding (A.T.7) and (A.T.8) yields
$4g+4\le8+s\le11$. Since $g\ge1$, we have
\[
           g=1,\qquad \deg\mathcal B=4,
                    q_0\in\{-2,-1,0\}.            \tag{A.T.9}
\]

We next prove the exact local obstruction required in
the equality cases. Let a smooth point of $D$ be a
branch point at which $\Delta$ has a simple zero.
The preceding nonzero-Hessian argument gives $a(0)=0$.
In coordinates $D=(x,y)=0$ with curve parameter $z$,
choose the holomorphic saturated Hessian factor as the
first normal coordinate. The quadratic jet is then
$a(z)x^2$. The radical cubic vanishes identically,
so a local reduced equation has the exact expansion
\[
 f_0=a(z)x^2+x^2(xA+yB)+xy^2b(z)+xy^3C+y^4E,
                                                        \tag{A.T.10}
\]
with holomorphic coefficients. The identity
$\Delta=ae-b^2/4$, $a(0)=0$ and the simple zero of
$\Delta$ imply $b(0)=0$, $a'(0)\ne0$ and $e(0)\ne0$.
Indeed an order at least two for $a$, together with
$b(0)=0$, would give an order at least two for $\Delta$;
and then $\Delta'(0)=a'(0)e(0)$ is nonzero.
Thus $b/a$ is holomorphic and $E(0)\ne0$.

Set $X=x+b(z)y^2/(2a(z))$. In this actual analytic
coordinate change the coefficient of $Xy^2$ cancels.
All terms created from $x^2(xA+yB)$ lie in
$(X^2)+(Xy^3)+(y^4)$. Therefore, collecting the
terms of $X$-degree zero and one, gives exactly
\[
 f_0=X^2A_1(X,y,z)+Xy^3B_1(y,z)+y^4C_1(y,z),
 \qquad C_1(0)\ne0.
\]
Choose a fourth root of $C_1$ and define
\[
 Y=yC_1^{1/4}+
                  \frac{XB_1}{4C_1^{3/4}}.
                                                        \tag{A.T.11}
\]
The map $(X,y,z)\mapsto(X,Y,z)$ is locally
biholomorphic, and direct expansion gives
$f_0-Y^4\in(X^2)$. Hence $f_0=Y^4+X^2H(X,Y,z)$.
Its quadratic jet along $D$ is still $a(z)X^2$,
so $H(0,0,z)=a(z)$ and $\partial H/\partial z(0)\ne0$.
The further coordinate $Z=H$ gives exactly
\[
                         f_0=Y^4+X^2Z.             \tag{A.T.12}
\]
To apply this to the actual map, first take a local
$m$th root of its unit. This multiplies the reduced
equation by a unit, which changes neither zero order
of $a$ nor of $\Delta$. The same coordinate construction
then makes the full map $(Y^4+X^2Z)^m$.
Statement A.5.15 excludes it. Consequently every
smooth branch point has order at least three in
$\operatorname{div}(\sigma)$.

We also record precisely when the conductor denominator
can be removed at a singular point $x\in D$.
If the intrinsic section $a$ is nonzero at every
normalization point above $x$, choose an actual
constant linear combination of the two lci conormal
generators which is nonzero under the normalized
quotient at all those finitely many points. The
corresponding diagonal entry $a_0$ of the symmetric
Hessian matrix belongs to $\mathcal O_{D,x}$ and
has nonzero value on every normalization branch.
It is therefore a unit in $\mathcal O_{D,x}$.
The other normalized factor coefficient is the
actual quotient $h=q_{XY}/a_0\in\mathcal O_{D,x}$.
Thus square completion and monic polynomial division
in the proof of (A.J.7)--(A.J.8) both take place
over $\mathcal O_{D,x}$ itself. The regular-sequence
graded pieces inject into their normalization
pullbacks, so their cancellations already descend
on $D$. We may repeat that proof with $d=1$.
It gives a holomorphic section $\Delta$ of $L^2$
above $x$, without a conductor pole. At a branch
point above $x$, this proves the sharper inequality
\[
 \operatorname{ord}\sigma\ge
                        4\operatorname{ord}\Xi+1.
                                                        \tag{A.T.13}
\]

We now exhaust the three integer cases in (A.T.9).
If $q_0=-2$, the divisor of $\sigma$ has degree four
and has four odd-multiplicity points, so each is
simple. At most three lie above singularities of
$D$, by (A.T.4). The remaining smooth branch point
contradicts (A.T.12).
If $q_0=0$, the section $a$ has no zeros and all
four branch points would have to lie above the
singular locus; this contradicts $s\le3$.

Finally let $q_0=-1$, so that $\deg\operatorname{div}(a)=2$
and $d_\Delta=6$. If $k$ of the four branch points
lie above $\operatorname{Sing}D$, then the local
obstruction and parity give
\[
               6\ge k+3(4-k)=12-2k.
\]
Hence $k=s=3$, and equality holds: the three singular
preimages have order one in $\operatorname{div}(\sigma)$
and the one smooth branch point has order three.
The relations $s\le\delta+\beta\le3$ now force
$\delta=2$, $\beta=1$ and two singular points of
delta invariant one. The exact delta-one
classification in A.5.14 identifies them as one
ordinary node and one cusp. Their normalization
conductor orders are respectively $(1,1)$ and $2$;
these follow directly from the displayed node and
cusp rings in that proof.

At the cusp preimage, if $a$ were nonzero,
(A.T.13) would make the order of $\sigma$ at least
$4\cdot2+1=9$, exceeding its total degree six.
So $a$ vanishes there. It also vanishes at the
smooth branch point. These two distinct points
exhaust the degree-two divisor of $a$. The section
is consequently nonzero at both node preimages.
Apply (A.T.13) at the node: each of its two branch
points has order at least $4\cdot1+1=5$ in
$\operatorname{div}(\sigma)$, again exceeding the
total degree six. This excludes every rational lci image.

We finish by treating all non-lci image germs permitted
by $\delta\le2$ and $\beta\le1$, with explicit analytic
rings and transverse coordinate maps. For a unibranch
reduced curve ring $A\subset\C\{t\}$ of colength at
most two, an element of order one makes $A=\C\{t\}$.
Otherwise its least positive valuation is two or three:
a least valuation at least four leaves three linearly
independent missing coefficients. If it is three and
the colength is two, the inclusion
$A\subset\C+t^3\C\{t\}$ is equality, since the latter
already has colength two. If it is two, choose a
parameter with $t^2\in A$ by taking a square root of
the leading unit. The even part of every convergent
series belongs to $\C\{t^2\}\subset A$. Subtracting
it shows that the least odd element can be written
$t^{2k+1}v(t^2)$; division by its unit puts
$t^{2k+1}$ in $A$. Consequently
$A=\C\{t^2\}\oplus t^{2k+1}\C\{t^2\}$, of colength
$k\le2$. These are the planar cusp and the planar
$A_4$ curve, besides the smooth case.

For two branches, write their image rings as $A_1,A_2$.
The two projections of $A\subset A_1\oplus A_2$ are
surjective. Their kernels define a common finite
quotient $Q_0$ and an exact fibre-product description
$A=A_1\times_{Q_0}A_2$. Thus
$\delta=\delta_1+\delta_2+\operatorname{length}Q_0$.
If both branches are smooth, the quotient is
$\C\{t\}/(t^k)$ with $k\le2$; after changing the
second parameter, the two maps identify the same
$k$-jet. The map
\[
 x\longmapsto(t,s),\qquad
 y\longmapsto(t^k,-s^k)
\]
identifies this fibre product with
$\C\{x,y\}/(y^2-x^{2k})$: write a matching pair as
the sum of its average and $y$ times its difference
divided by $2x^k$. These are the node and tacnode.
If one branch is singular, the delta budget forces
one cusp, one smooth branch, and $Q_0=\C$.
Three branches are excluded by $\beta\le1$.

An abstract plane curve ring in the given smooth
threefold is lci by the actual transverse-coordinate
construction in A.5.14. Thus the only remaining
non-lci analytic rings, necessarily consuming the
whole delta budget at a single point, are
\[
\begin{array}{ll}
 A=\C\{t^3,t^4,t^5\},
 &I=(y^2-xz,\ yz-x^3,\ z^2-x^2y),\\[2pt]
 A=\C\{t^2,t^3\}\times_{\C}\C\{s\},
 &I=(xz,\ yz,\ y^2-x^3).
\end{array}                                             \tag{A.T.14}
\]
These are normal forms of the actual embeddings.
Their three displayed coordinate functions form a
basis of the cotangent space of each ring. Lift them
to $\mathcal O_{Y,x}$; their differentials are linearly
independent because the embedding dimension is three.
They are therefore ambient coordinates by the inverse
function theorem. The displayed ideals are their
kernels, obtained respectively by the parametrizations
$(t^3,t^4,t^5)$ and the union of $(t^2,t^3,0)$ with
$(0,0,s)$.

We show that the local reduced equation $f_0$ actually
lies in $I^2$ in each case. In the first case put
\[
 g_1=y^2-xz,\quad g_2=yz-x^3,\quad
 g_3=z^2-x^2y,\quad
 F=x^5+xy^3+z^3-3x^2yz.
\]
The exact identities
\[
 xF=g_2^2-g_1g_3,\qquad
 yF=xg_1^2+g_2g_3,\qquad
 zF=g_3^2+xg_1g_2                              \tag{A.T.15}
\]
show that the maximal ideal annihilates the class of
$F$ modulo $I^2$. Moreover
$I^{(2)}=I^2+(F)$ in the convergent local ring.
Here is a finite-module verification of that equality.
Let $J=I^2+(F)$. The quotient $\C\{x,y,z\}/J$ is
finite over the discrete valuation ring $\C\{x\}$,
since its reduction is the finite parametrized curve.
Its generic rank is nine: the curve has degree three
over the $x$ coordinate, and the square of its smooth
codimension-two ideal has length three along each
generic normal slice. Modulo $x$, the quotient has
the nine-element basis
\[
 1,y,y^2,y^3,\ z,yz,y^2z,\ z^2,yz^2,
\]
because its ideal is $(y^4,y^3z,y^2z^2,z^3)$.
A finite module over a discrete valuation ring is
the sum of its free and torsion parts; equality of
generic rank and residue-space dimension here makes
the torsion part zero. Thus $J$ is $x$-saturated.
By (A.T.15), $J$ and $I^2$ agree after inverting $x$.
The curve is lci away from the origin, so its symbolic
square differs from its ordinary square only by a
module supported there. Hence
$I^{(2)}=I^2:x^\infty=J$.

On its normalization let
$a_1=dg_1$ and $b_1=dg_2+t\,dg_1$.
They form a basis of the normalized conormal image;
$dg_3=-2t^2a_1+t b_1$.
The quadratic term of $F$ in this basis is
\[
 q_F=\frac3t a_1^2-\frac3{t^2}a_1b_1
                           +\frac1{t^3}b_1^2,
 \qquad \det(q_F)=\frac3{4t^4}.                  \tag{A.T.16}
\]
Every equation in $I^2$ has a holomorphic quadratic
form in this same basis. If $f_0=cF+I^2$ had
$c(0)\ne0$, the determinant of its quadratic form
would have the noncancellable leading pole
$3c(0)^2/(4t^4)$: the cross terms with a holomorphic
form have pole at most three. It could not have
generic rank one. Thus $c(0)=0$, and (A.T.15) gives
$f_0\in I^2$.

For the second ring put $G=y^2-x^3$ and
$P=(z,G)$, $Q_1=(x,y)$. Then $I=P\cap Q_1$ and
\[
 I^{(2)}=P^2\cap Q_1^2=I^2+(zG).                \tag{A.T.17}
\]
Indeed $P^2=(z^2,zG,G^2)$, while $zG,G^2$ already
belong to $Q_1^2$. Multiplication by $z^2$ is
injective modulo $Q_1^2$, so membership of
$z^2A+zGB+G^2C$ in $Q_1^2$ forces $A\in Q_1^2$.
This gives the asserted generating ideal directly.
Furthermore
$x(zG)=(xz)G$, $y(zG)=(yz)G$, and
$z(zG)=(yz)^2-x(xz)^2$ all lie in $I^2$.
On the cusp branch use $a_1=d(xz)=t^2dz$ and
$b_1=dG=-3t^4dx+2t^3dy$.
The quadratic form of $zG$ is $t^{-2}a_1b_1$,
of determinant $-1/(4t^4)$. The same rank-one
argument forces its coefficient to vanish at the
origin, and the three identities again give
$f_0\in I^2$.

For clarity, the bundle used at these non-lci points
is now defined by its actual morphism. Put
\[
 K=\operatorname{im}\bigl(\nu^*(I/I^2)
                         \longrightarrow\nu^*\Omega_Y^1\bigr),
 \qquad N=K^*.
\]
The coherent image $K$ is locally free of rank two
on the smooth curve $\Sigma$. Let $R_{\rm ram}$ be
the zero divisor of $d\nu$, and let $K_{\rm sat}$
be the kernel of the surjection
$\nu^*\Omega_Y^1\to\Omega_\Sigma^1(-R_{\rm ram})$.
The torsion quotient $K_{\rm sat}/K$ defines a divisor
$J_{\rm jac}$. Taking determinants in these actual
maps gives
\[
 \deg N=-2+\deg(J_{\rm jac}-R_{\rm ram}).         \tag{A.T.18}
\]
For $(t^3,t^4,t^5)$, the wedge $a_1\wedge b_1$
has coefficients $t^8(3,4t,5t^2)$, while $d\nu$
has common factor $t^2$. Thus the correction is
$8-2=6$. For the cusp branch in the second case,
the wedge is $-t^5(2,3t,0)$ and the tangent common
factor is $t$, giving correction four. On its
smooth branch the wedge is $s^2(dx\wedge dy)$,
giving correction two. In both cases (A.T.18)
therefore gives $\deg N=4$.

Because $f_0\in I^2$, its quadratic form is a
holomorphic section of $\operatorname{Sym}^2K\otimes M$,
independent of its expression as a sum of products
of ideal generators. This follows either from its
actual Hessian or from injectivity of
$\operatorname{Sym}^2K\to\nu^*\operatorname{Sym}^2\Omega_Y^1$.
Its saturated radical and quotient again give line
bundles $R,Q$, with $q_0=\deg Q\le0$ and
$\deg R=4-q_0$. The generic quartic invariant
is meromorphic in $L^2$, where
$L=M\otimes Q^{-1}\otimes R^{-2}$; now
\[
                        \deg L^2=-16+2q_0.        \tag{A.T.19}
\]

We bound its poles using actual transverse affine
coordinate maps, with the curve parameter held
fixed during every normal differentiation. For the
space cusp use, for $t\ne0$,
\[
 (u,v,t)\longmapsto
 \left(t^3,\ t^4+\frac{v}{3t^5},\
                   t^5-\frac{u}{t^3}+\frac{2v}{3t^4}\right).
                                                        \tag{A.T.20}
\]
This is locally biholomorphic and its two transverse
derivatives are the dual frame to $a_1,b_1$.
Each coefficient of a derivative has pole at most
five. The first normal derivatives of all three
$g_i$ are holomorphic, since
$dg_1=a_1$, $dg_2=b_1-ta_1$ and
$dg_3=-2t^2a_1+t b_1$.
Second normal derivatives have pole at most ten.
For $f_0=\sum c_{ij}g_i g_j$, the cubic Taylor
coefficients consequently have pole at most ten:
the only surviving terms use a first and a second
derivative of the $g_i$, or a first derivative of
$c_{ij}$ and two first derivatives of the $g_i$.
Every fourth derivative of a holomorphic ambient
function in (A.T.20) has pole at most twenty.

Choose a holomorphic change of the $N$ frame adapted
to its saturated radical and quotient. It does not
increase these bounds. The radical cubic is zero
because the generic transversal is $A_3$.
Writing the relevant Taylor terms as
$a\ell^2+b\ell r^2+er^4$ therefore gives a
holomorphic $a$, a mixed cubic $b$ of pole at most
ten, and $e$ of pole at most twenty. The invariant
$\Delta=ae-b^2/4$ has pole at most twenty.
These are the same intrinsic jets on overlaps:
to weighted order four, where $\ell$ has weight
two and $r$ has weight one, an adapted change is
$\ell'=\alpha\ell+u r^2$, $r'=\beta r$.
Direct substitution multiplies $\Delta$ by
$\alpha^{-2}\beta^{-4}$; changes of the defining
section supply the factor $M^2$. Higher weighted
terms and changes of the curve parameter do not
alter these coefficients. This verifies the
transition law for the line $L^2$ in (A.T.19).

There is only one singular preimage $p$ in this case.
Thus $\Delta$ is a holomorphic section of the square
line $L^2(20p)$, whose degree is $4+2q_0$.
As in (A.T.7), the even pole correction preserves
branch parity, so
\[
 \deg\mathcal B\le4+2q_0,
 \qquad \deg\mathcal B\le-2q_0+1.
\]
The second bound again uses A.5.17 away from that
single singular image preimage. Adding them gives
$2\deg\mathcal B\le5$, contradicting
$\deg\mathcal B=2g+2\ge4$.

For the cusp-plus-line germ, the corresponding
actual affine coordinate maps are
\[
 (u,v,t)\longmapsto
       \left(t^2,t^3+\frac{v}{2t^3},\frac{u}{t^2}\right),
 \qquad
 (u,v,s)\longmapsto\left(\frac us,\frac vs,s\right).
                                                        \tag{A.T.21}
\]
Their first generator derivatives are holomorphic
in the respective normalized conormal frames.
The same product differentiation gives cubic pole
bounds six and two, and the fourth derivative bounds
are twelve and four. Thus $\Delta$ has poles bounded
by $12p+4q$ at the two normalization preimages.
The line $L^2(12p+4q)$ is a square and has degree
$2q_0\le0$ by (A.T.19). A nonzero section has no
positive-degree odd divisor, whereas this double
cover has at least four branch points. This is the
last contradiction. All non-lci cases have now
been excluded, and the positive-genus case at the
start gives the asserted elliptic curves and constant
transversal type.

We finally prove the normalization and Euler assertions,
keeping track of any unnecessary blowups in the original
choice of resolution. Factor $\pi$ through the minimal
resolution $\pi_{\min}:Z_{\min}\to T$, with the resulting
sequence of exceptional-curve blowdowns
$q:Z\to Z_{\min}$. Every curve contracted by $q$ is
vertical for $p$, by the stated assumption on $\pi$.
The equality of structure-sheaf direct images for each
blowdown, in local coordinates on $B$, gives a descended
ruling $p_{\min}:Z_{\min}\to B$ with
$p=p_{\min}q$. The image $C_{\min}$ is still a section,
and the pushed-forward effective conductor is
$D_{\min}=2C_{\min}+V_{\min}$ with $V_{\min}$ vertical.
The normalization-canonical comparison gives
$K_{Z_{\min}}+D_{\min}=\pi_{\min}^*\eta^*\omega_S$
as line bundles, and its rational Chern class is zero.

Statement C2d gives a smooth normalization near the
inverse image of the elliptic curve $D$, with normal
line of degree zero on its unramified double cover.
The minimal resolution is an isomorphism there, so
$C_{\min}^2=0$. No other conductor or exceptional
curve meets it. Hence $V_{\min}\cdot C_{\min}=0$.
No ruling $(-1)$-curve can meet $C_{\min}$: one in
$V_{\min}$ is disjoint from it, while one outside
$V_{\min}$ would have
\[
 (K_{Z_{\min}}+2C_{\min}+V_{\min})\cdot E
       =-1+2C_{\min}\cdot E+V_{\min}\cdot E\ge1,
\]
contrary to the vanishing Chern class.

Contract ruling $(-1)$-curves to a minimal ruled bundle.
Every contraction therefore avoids the section. For
one blowdown $\sigma:W\to W'$ let $D'=\sigma_*D_W$,
let $a$ be the coefficient of its exceptional curve,
and let $d$ be the multiplicity of $D'$ at the centre.
The actual divisor comparison is
\[
 K_W+D_W=\sigma^*(K_{W'}+D')+(1+a-d)E.
                                                        \tag{A.T.22}
\]
Its intersection with $E$ gives $a=d-1$, so the
log-canonical line and its zero rational Chern class
descend. The section remains disjoint from the
effective vertical part at every step. On a minimal
ruled bundle every nonzero effective vertical divisor
is a positive sum of fibres and meets the section.
The final vertical part is consequently zero. If
there had been a blowup in the reverse sequence,
its first centre would be away from this section,
and (A.T.22) would give coefficient $-1$ on its
exceptional curve. Its strict transform retains that
negative coefficient under subsequent blowups, a
contradiction to effectiveness. Thus no ruling
blowdowns were needed: $Z_{\min}$ is already a
minimal ruled bundle and $D_{\min}=2C_{\min}$.

All exceptional curves of $\pi_{\min}$ would be
vertical. A minimal ruled bundle has irreducible
fibres of square zero, which cannot form the
negative-definite exceptional divisor of a surface
resolution. Hence $\pi_{\min}$ has no exceptional
curve and is an isomorphism. This proves the assertion
about $T$. Finally $\eta$ is an isomorphism away from
its conductor section $C_T$ and the elliptic image
$D$. Additivity of the topological Euler number on
these actual decompositions gives
\[
 e(S)=e(T)-e(C_T)+e(D)
      =e(\PP^1)e(B)-0+0=0.                       \tag{A.T.23}
\]
The full multiple fibre has the same underlying
topological space and therefore the same Euler number.
\end{proof}

\paragraph{Connected double covers in the two non-lci cases.}
The non-lci calculation also excludes a connected
degree-two conductor cover when its covering curve
is rational, for either of the two image rings in
(A.T.14). In the cusp-plus-line case the cleared
quartic has degree $2q_0\le0$, so it has no branch
divisor at all. In the space-cusp case a rational
double cover has branch degree two. The already
proved inequalities give
\[
                2\le4+2q_0,\qquad 2\le-2q_0+1,
\]
so the integer $q_0$ is $-1$. The cleared quartic
then has degree two and exactly two simple branch
points. Only one point lies above the singular
image germ. The other is a smooth simple branch
point, excluded by the actual coordinate reduction
(A.T.10)--(A.T.12). Thus these two non-lci images
cannot support a connected double cover of any
nonnegative genus in a torus degeneration.

The integer branch budgets, the discriminant square
weights and the coordinate divisibility used in
(A.T.10)--(A.T.12) are checked in
\path{code/E_tacnode_bounds.py}; the non-lci charts,
symbolic generators and cubic/quartic pole orders
are checked in \path{code/E_low_delta_tacnode.py}.
The topological
quotient (A.T.2) targets the actual singular curve
$D/J$; replacing it by its normalization would
lose the free-group obstruction in (A.T.4).

\subsection{The image curve of an ordinary-triple conductor}

\begin{statement}{A.5.19}{Proved}{Ordinary-triple images have arithmetic genus at most one}
Let $f:Y\to C_0$ be a proper holomorphic map from a smooth
compact complex threefold with $b_2(Y)=0$, with connected
fibres and smooth general fibre a complex two-torus.
Let $mS$ be a full fibre with irreducible reduced support,
and let $D\subset\operatorname{Sing}S$ be an irreducible
curve with generic ordinary-triple plane transversal.
Then $D$ is a local complete intersection in $Y$ and
$p_a(D)\le1$. The possibilities for its underlying curve
are a smooth rational curve, a rational curve with one
ordinary node or cusp, and a smooth elliptic curve.
In the elliptic case the transversal type is constant.

When $D$ is singular rational, every component of the
normalized conductor curve dominating $D$ is rational,
and its map to $D^\nu\simeq\PP^1$ has degree one.
Consequently, in the irrational ruled setup of A.5.18
with the generic transversal replaced by an ordinary
triple point, the image of the noncontracted horizontal
section is either smooth rational or smooth elliptic.
In the elliptic case the horizontal cover is unramified
of degree three, the normalization is smooth minimal
elliptic ruled, and the fibre has Euler number zero.
In the rational case the horizontal cover has degree
two or three; these cases are not excluded by this
statement.
\end{statement}
\begin{proof}
By A.5.14, $p_a(D)\le2$. If its normalization has
positive genus, that statement and A.5.5 already give
a smooth elliptic curve of constant transverse type.
Suppose its normalization $\nu:\Sigma\to D$ is
$\PP^1$, with total delta invariant $\delta\le2$.

First let $D$ be lci. Put
$N=(\nu^*(I_D/I_D^2))^*$ and
$M=\nu^*(\mathcal O_Y(S)|_D)$.
The quadratic normal jet vanishes generically, so
its vanishing descends on the reduced curve through
the locally free graded piece $I_D^2/I_D^3$.
Thus the initial cubic is an actual section of
$\operatorname{Sym}^3N^*\otimes M$.
Adjunction and the curve conductor give
$\deg N=-2+2\delta$, while $\deg M=0$.
Its nonzero binary-cubic discriminant is a section
of $M^4\otimes(\det N)^{-6}$, of degree
\[
                         12-12\delta.             \tag{A.U.1}
\]
Hence $\delta\le1$.

We exclude every non-lci germ of delta at most two
by its symbolic cube. The one- and two-branch analytic
classification is proved in (A.T.14). There is one
additional possibility when three branches occur:
delta at least two is forced by identifying their
constants, so equality makes the ring exactly
\[
 \{(a,b,c)\in\C\{t_1\}\oplus\C\{t_2\}\oplus\C\{t_3\}:
                      a(0)=b(0)=c(0)\}.
\]
Its actual embedding in the smooth threefold is the
union of the three coordinate axes, with ideal
$I=(xy,xz,yz)$. Indeed the three branch parameters
give a basis of the cotangent space, so their lifts
are ambient coordinates as in A.5.18.
Thus there are exactly three non-lci cases to consider;
each consumes the whole delta budget at one point.

For the space cusp retain $I=(g_1,g_2,g_3)$ and
$F$ from (A.T.15). Its exact symbolic-cube identity is
\[
                         I^{(3)}=I^3+FI.           \tag{A.U.2}
\]
Here is a verification in the convergent local ring.
The right-hand side agrees with $I^3$ after inverting
$x$, by (A.T.15). Its quotient is finite over
$\C\{x\}$, of generic rank $3\cdot6=18$: a square
normal plane has six monomials of total degree less
than three. Modulo $x$, its ideal in $\C\{y,z\}$ is
\[
 (y,z)^6+(y^2z^3,yz^4,z^5).
\]
The fifteen monomials of total degree at most four,
together with $y^5,y^4z,y^3z^2$, give residue dimension
eighteen. The same discrete-valuation-module argument
used for (A.T.15) makes the quotient $x$-torsion-free.
It is therefore the saturation of $I^3$, which equals
the symbolic cube since the curve is lci off the origin.
This proves (A.U.2).

For the cusp-plus-line ideal $I=(xz,yz,G)$ with
$G=y^2-x^3$, one has
\[
                   I^{(3)}=I^3+(zG)I.             \tag{A.U.3}
\]
To verify this directly, write $P=(z,G)$ and
$Q=(x,y)$. Their powers are primary because these
prime ideals are complete intersections. The
symbolic cube is $P^3\cap Q^3$.
Write an element of $P^3$ as
$z^3A+z^2GB+zG^2C+G^3E$.
The last two terms are already in $Q^3$.
Reducing the other terms modulo $Q^3$ gives
$zA+GB=0$ there. The constant-in-$(x,y)$ coefficient
of $B$ must therefore be divisible by $z$; write
$B=zB'+B_Q$, $B_Q\in Q$.
It follows that $A+GB'\in Q^3$.
Consequently
\[
 P^3\cap Q^3=(z^3Q^3,z^2GQ,zG^2,G^3)
                        =I^3+(zG)I,
\]
as claimed.

For the three axes, membership of
$x^a y^b z^c$ in the symbolic cube means
$a+b,a+c,b+c\ge3$.
If one exponent is zero, the other two are at
least three and the monomial belongs to $I^3$.
If all are positive, either at least two are at
least two, in which case it belongs to $(xyz)I$,
or two are one, which violates one of the three
inequalities. Hence
\[
                         I^{(3)}=I^3+(xyz)I.       \tag{A.U.4}
\]

In all three cases the generic ordinary-triple
equation $f_0$ belongs to $I^{(3)}$: on the smooth
locus of each curve branch its order in the two
normal variables is at least three. Define the
bundle by its actual conormal image,
\[
 K=\operatorname{im}\bigl(\nu^*(I/I^2)
                           \to\nu^*\Omega_Y^1\bigr),
 \qquad N=K^*.
\]
For the first two cases (A.T.18) gives $\deg N=4$.
For the three axes, on the branch $(t,0,0)$ the
frame is $t\,dy,t\,dz$, with conormal determinant
order two and no tangent ramification. Each of the
three branches has the same correction, for total
six, so again $\deg N=-2+6=4$.

Use the actual affine maps (A.T.20)--(A.T.21).
For the three axes add
\[
              (u,v,t)\longmapsto(t,u/t,v/t)
                                                        \tag{A.U.5}
\]
on the first branch, and the two analogous permuted
maps on the other branches. In these charts the
first normal derivatives of the ideal generators
are holomorphic in the specified conormal frame.
An $I^3$ term consequently has holomorphic initial
cubic. The extra term in each of (A.U.2)--(A.U.4)
has initial cubic equal to the quadratic term of
its displayed symbolic-square generator, multiplied
by the first derivative of an ideal generator.

For the space cusp that quadratic term is (A.T.16),
so the cubic pole is at most three. For the
cusp-plus-line germ it is $t^{-2}uv$ on the cusp
and $s^{-1}v^2$ on the smooth branch, so the pole
bounds are two and one. For the three axes it is
$t^{-1}uv$ in (A.U.5), giving pole bound one on
each branch. Coefficients multiplying these terms
are holomorphic; their derivatives do not enter
the initial cubic because the lower normal jets
of the products vanish.

The initial cubic is intrinsic on the dense smooth
part of $D$, as all lower jets are zero. Its
transition under a normal-coordinate change is
therefore the actual cubic tensor in
$\operatorname{Sym}^3N^*\otimes M$; higher-order
coordinate terms make no contribution. The binary
discriminant is homogeneous of degree four in
these coefficients. It is a meromorphic section
of $M^4\otimes(\det N)^{-6}$, of degree $-24$,
with pole divisor of total degree at most twelve
in each of the three cases. After this exact pole
correction its line has negative degree, at most
$-12$. It cannot have the required nonzero
holomorphic section. This excludes all non-lci
cases and proves that $D$ is lci with $p_a(D)\le1$.
The list of underlying curves follows from the
delta-one classification in A.5.14.

We give the finite-cover comparison at a singular
rational image. Here $\delta=1$, so (A.U.1) has
degree zero and its nonzero section is nowhere zero
on $\Sigma=\PP^1$. The cubic defines a subscheme
$R_3\subset\PP(N)$ of tangent-root lines. Every
fibre consists of three distinct points. In a
projective coordinate chart, its derivative with
respect to the root coordinate is nonzero, so the
implicit function theorem makes
$R_3\to\Sigma$ a finite etale map of degree three.
There is no component in a whole projective fibre,
since its nonzero discriminant forbids an identically
zero cubic there.

Over the generic point of $D$, normalization of an
ordinary triple transversal separates exactly its
three tangent lines. This supplies an isomorphism
between the function-field algebra of $R_3$ and
that of the normalized conductor curves dominating
$D$, after their maps are factored through $\Sigma$.
Both are normal finite curves over $\Sigma$.
Taking the integral closure of $\mathcal O_\Sigma$
in this same algebra therefore identifies the two
finite covers everywhere. Every connected finite
etale cover of the simply connected Riemann sphere
has degree one. Thus each of these conductor
components is rational and maps isomorphically to
$\Sigma$, proving the second assertion.

In the irrational ruled setting the horizontal
section has positive genus. Its finite map to a
rational image normalization must have degree at
least two. The preceding comparison excludes a
singular rational image, and the generic triple
transversal bounds its degree above by three.
The remaining rational image is therefore smooth,
with degree two or three. If the image is elliptic,
no vertical rational curve can dominate it, so all
three branches belong to the horizontal section.
Statement C2d gives an unramified degree-three
cover, normal degree zero and constant type.
The exact minimal-resolution factorization and
blowdown argument (A.T.22) applies without change
to this coefficient-two section: its proof uses
only those properties. It makes the normalization
smooth minimal elliptic ruled with conductor twice
the section. The actual normalization Euler formula
(A.T.23) then gives zero.
\end{proof}

The singular rational possibilities cannot be
removed by a blanket local spherical-cycle argument.
For example
\[
                     f=x(yz)(x-yz)
                                                        \tag{A.U.6}
\]
has ordinary-triple transversals along the nodal curve
$D=(x,yz)=0$ and an additional ordinary-double
critical line $y=z=0$. For $t\ne0$, put
$B_t=\{(x,v):xv(x-v)=t\}$. The exact maps
\[
 B_t\times\C^*\longrightarrow f^{-1}(t),
 \quad (x,v,y)\longmapsto(x,y,v/y),\qquad
 (x,y,z)\longmapsto(x,yz,y)
                                                        \tag{A.U.7}
\]
are inverse holomorphic maps. Nonzero $t$ forces
$v,y$ and $z$ to be nonzero. The smooth affine
curve $B_t$ and $\C^*$ are aspherical, so this
regular affine fibre has no spherical homotopy
class. This is a local model, not a compact
$b_2=0$ example. A deformation smoothing $yz$
changes the additional critical line's transverse
type on the boundary and does not satisfy the
comparison used in A.5.17.

Exact ideal, normal-coordinate and discriminant
calculations are retained in
\path{code/E_triple_images.py} and
\path{code/E_triple_images.out}; the independently
computed symbolic saturations are also in
\path{code/A_low_delta_cusp.py}.

\subsection{Completion of the noncontracted horizontal-cusp case}

\begin{statement}{A.5.20}{Proved}{A noncontracted horizontal cusp has zero Euler number}
Let $f:Y\to C_0$ be a proper holomorphic map from a smooth
compact complex threefold with $b_2(Y)=0$, with connected
fibres and smooth general fibre a complex two-torus.
Let $mS$ be a full fibre with irreducible reduced support,
$\eta:T\to S$ its normalization and $\pi:Z\to T$ an
irrational ruled resolution with ruling $p:Z\to B$.
Suppose every exceptional curve of $\pi$ is vertical
and the resolution conductor is $2C+V$, where $C$
is a noncontracted section and $V$ is vertical.
If the generic transversal at $D=\eta\pi(C)$ is an
ordinary cusp, then $D$ and $C$ are smooth elliptic
curves and $C\to D$ is an isomorphism. The cusp type
is constant along $D$. The normalization $T$ is a
smooth minimal ruled surface over this elliptic curve,
its conductor is twice a section, and $e(S)=0$.
This permits an initially singular normalization;
smoothness is part of the conclusion.
\end{statement}
\begin{proof}
A transverse cusp has exactly one normalization
branch. The map from the horizontal section to the
normalization of its image consequently has degree
one. Both curves are smooth, so this finite map
identifies $C\simeq B$ with $D^\nu$. Its genus is
positive. Statement A.5.14, including its actual
delta-one lci comparison, and A.5.5 therefore make
$D$ itself smooth elliptic and its cusp transversal
constant at every point. Thus $C\to D$ is an
isomorphism, without any assumption about the
ambient Picard pullback.

For the exact local normalization $x=u^3,y=u^2$,
let $C_T\subset T$ be the image of $C$. The
normalization is smooth near this curve. Its normal
line $P=N_{C_T/T}$ has $P^2=R$ and $P^3=Q$
in the degree-zero radical and Hessian quotient
lines of C2d. Hence $\deg P=0$. The normalization
is smooth near $C$, and no other conductor curve
meets it. Factor $\pi$ through its minimal
resolution, descending the ruling across its
vertical exceptional blowdowns as in A.5.18.
On that minimal resolution the section is isomorphic
to $C_T$ in its unchanged smooth neighbourhood. It
therefore has self-intersection zero and is disjoint from
the effective vertical conductor. The canonical
comparison is the actual line-bundle identity
$K+2C+V=(\eta\pi)^*\omega_S$, of zero rational
Chern class.

Apply (A.T.22) to each ruling blowdown. Its
exceptional coefficient is exactly the multiplicity
at its centre minus one. No ruling $(-1)$-curve
can meet $C$, since the intersection with $K+2C+V$
would then be positive. On the minimal ruled
bundle the vertical conductor is zero because it
is disjoint from a section. A first reverse blowup
away from that section would have coefficient
$-1$, so there were no ruling blowups. All
resolution-exceptional curves are vertical, and
the square-zero fibres of the resulting bundle
cannot be exceptional. This identifies the minimal
resolution with $T$ itself and proves the asserted
smoothness and minimality.

Finally the normalization is an isomorphism off
$C$ and $D$. The exact Euler comparison is
$e(S)=e(T)-e(C)+e(D)$; all three terms vanish,
since $T$ is ruled over an elliptic curve and
$C,D$ are elliptic. Thus $e(S)=0$, also for the
underlying topological space of the full fibre $mS$.
\end{proof}

\subsection{Binary products over a nodal plane function}

\begin{statement}{A.5.21}{Proved}{A homogeneous binary-product grid is excluded}
Let $f:Y\to C_0$ be a proper holomorphic map from a smooth
complex threefold with smooth nearby fibres complex
two-tori. Let $P(X,V)$ be a homogeneous polynomial of
degree $d\ge2$, a product of distinct complex linear
forms, and suppose $P(1,0)\ne0$.
There is no point at which the actual fibre map has
the local equation
\[
             f=u(x,y,z)\,P(x,yz)^m,
             \qquad m\ge1,                         \tag{A.BP.1}
\]
with $u$ a holomorphic unit. In particular this excludes
$\prod_{i=1}^d(x-a_i yz)$ for distinct constants $a_i$,
and every positive full-fibre multiplicity of that germ.
No Betti-number assumption on $Y$ is required.

The same exclusion holds for the variable-coefficient
equation $f=u h^m$ with
\[
 h=\sum_{i=0}^d a_i(x,y,z)x^{d-i}(yz)^i,
 \qquad P_0(X,V)=\sum_{i=0}^d a_i(0)X^{d-i}V^i,
                                                        \tag{A.BP.4}
\]
provided $P_0$ has $d$ distinct projective roots and
$P_0(1,0)\ne0$. The comparison with its constant
initial form is the explicit analytic family below.
\end{statement}
\begin{proof}
Put $h=P(x,yz)$. It is weighted homogeneous of degree
$2d$ for the positive weights $(2,1,1)$.
On a small simply connected coordinate neighbourhood
take a holomorphic logarithm of the unit and set
\[
 \Phi(x,y,z)=\bigl(xu^{1/(md)},
                  yu^{1/(2md)},zu^{1/(2md)}\bigr).
                                                        \tag{A.BP.2}
\]
Its derivative at the origin is invertible, and direct
homogeneity gives $h\Phi=u^{1/m}h$. Thus
$f=(h\Phi)^m$ exactly. It suffices to construct a
compact spherical cycle of nonzero square in $h=1$,
because weighted homotheties then put it inside an
arbitrarily small actual fibre of (A.BP.1).

The curve $B=\{P(x,v)=1\}\subset\C^2$ is smooth:
Euler's identity $xP_x+vP_v=dP$ precludes a common
zero of its two partial derivatives on that level.
It is also connected. The map
\[
 B\longrightarrow\PP^1\setminus\{P=0\},
                       \quad(x,v)\longmapsto[x:v]
\]
is an etale cover of degree $d$. Locally a projective
representative has nonzero $P$-value, and its $d$
rescalings with $P=1$ give the sheets. Around any
one of the simple projective zeros of $P$, that
value has a simple zero, so its inverse $d$th root
has monodromy a full $d$-cycle. This monodromy is
transitive on the sheets, proving connectedness.

The function $v$ on $B$ has exactly $d$ distinct
simple zeros, since
$P(x,0)=P(1,0)x^d$ and $P_x(x,0)=d/x$ at those
points. Choose an embedded arc $\gamma$ in $B$
joining two of its zeros and avoiding all other
zeros. On the interior choose a continuous argument
of $v(\gamma(t))$ and two functions $A(t),B_0(t)$
with $A B_0=v(\gamma(t))$ and
$|A|=|B_0|=|v(\gamma(t))|^{1/2}$.
The map
\[
 (t,\theta)\longmapsto
   \bigl(x(\gamma(t)),A(t)e^{\mathrm i\theta},
                        B_0(t)e^{-\mathrm i\theta}\bigr)
                                                        \tag{A.BP.3}
\]
lies in $h=1$ and collapses both boundary circles
of $[0,1]\times S^1$ to points. It is therefore a
spherical cycle, with embedded interior circles.

A parallel-arc isotopy with the same endpoints gives
a homologous spherical cycle whose support meets
the first only at those endpoints. At each endpoint,
$P_x\ne0$, so $(y,z)$ are actual holomorphic
coordinates on $h=1$. The links of the two cycles
are parallel curves of slope $(1,-1)$ on a Clifford
torus in the local three-sphere. Their linking
number is $-1$. As in the explicit relative-chain
comparison in A.5.15, that number is the local
intersection of the two coned cycles. Both
orientations reverse at the second endpoint,
so its intersection sign is unchanged. The
spherical class consequently has square $-2$.

The weighted homothety
$(x,y,z)\mapsto(\lambda^2x,\lambda y,\lambda z)$
maps $h=1$ biholomorphically to $h=\lambda^{2d}$
and shrinks the compact cycles and their isotopy
to the origin. Together with (A.BP.2), this puts
them inside an actual smooth global fibre of $f$.
That fibre is a complex two-torus and has zero
$\pi_2$, so the spherical class has zero homology
image. The oriented codimension-zero embedding
preserves the intersection of the compact cycles,
giving both zero and $-2$ for its square. This
contradiction proves the statement.

We now prove the variable-coefficient assertion,
including the boundary comparison. Absorb a local
$m$th root of the unit into $h$; this preserves the
form (A.BP.4) and its two initial-form conditions.
For a complex parameter $s$ set
\[
 h_s=\sum_{i=0}^d a_i(s^2x,sy,sz)x^{d-i}(yz)^i.
                                                        \tag{A.BP.5}
\]
This is holomorphic also at $s=0$, with
$h_0=P_0(x,yz)$. For $s\ne0$ there is the exact
identity
\[
       h(s^2x,sy,sz)=s^{2d}h_s(x,y,z).           \tag{A.BP.6}
\]
The leading coefficient $a_0$ is a unit. The implicit
function theorem applied to its degree-$d$ root
polynomial gives distinct holomorphic functions
$b_i(s,x,y,z)$ and the exact factorization
\[
       h_s=a_0(s^2x,sy,sz)
                     \prod_{i=1}^d(x-b_i(s,x,y,z)yz)
                                                        \tag{A.BP.7}
\]
on one small parameter and coordinate neighbourhood.
Each factor defines a smooth surface, since its
$x$ derivative is a unit. For $i\ne j$ the ideal
generated by the two factors equals $(x,yz)$:
their difference is the unit $b_j-b_i$ times $yz$.
Thus the zero level has precisely the same two
critical curves $D=(x,yz)=0$ for every $s$.

Choose a small closed ball $K$ in the coordinate
chart such that the two axes of $D$ meet its boundary
transversely and the zero-level surfaces away from
those axes are transverse to the boundary. At
$s=0$ this follows directly for the smooth graph
surfaces $x=b_i(0)yz$ and sufficiently small radii.
It persists for small $s$ by compactness away from
neighbourhoods of the two boundary links.
Near either link, use normal coordinates $(U,V)$
for the smooth critical curve and a parameter $w$
along it. The normal initial term of $h_s$ is a
degree-$d$ binary form with distinct roots.
On the compact parameterized link its normal
gradient therefore satisfies a uniform bound
\[
 \|\partial_{U,V}h_s\|\ge c\|(U,V)\|^{d-1},
 \qquad |\partial_w h_s|\le C\|(U,V)\|^d
                                                        \tag{A.BP.8}
\]
after the neighbourhood is shrunk. The first bound
follows because the gradient of a homogeneous
reduced binary form has no zero on the unit sphere;
higher normal terms have one additional order.
The second follows from $h_s\in(U,V)^d$.

Transversality of the critical curve to $\partial K$
allows a boundary coordinate with one real component
of $w$ solved as a smooth function of $(U,V)$ and
the remaining component. The correction to normal
differentiation tangent to this boundary is a
bounded multiple of $\partial_w h_s$.
By (A.BP.8) it is smaller by one normal order than
the surjective complex normal derivative. Thus
$h_s|_{\partial K}$ is a real submersion at every
nearby point with $h_s\ne0$, uniformly in $s$.
Away from the two link neighbourhoods, compactness
and the preceding zero-level transversality give
the same conclusion for all sufficiently small
nonzero values.

Remove the proper critical-value discriminant from
the product of the parameter disc and a sufficiently
small punctured value disc. The critical locus over
that domain cannot meet the boundary, so its
projection is proper and its image is complex
analytic. Its complement is connected. The level
sets in $K$ form a proper submersion with boundary
there, and Ehresmann transport identifies their
compact fibres by oriented diffeomorphisms.
Choose a small sphere cycle of the already proved
model $h_0=P_0(x,yz)$, and transport it to a fibre
of $h_{s_0}$ for $s_0\ne0$. Equation (A.BP.6)
then embeds that fibre into the actual fibre
of $h^m$ with value $s_0^{2dm}\epsilon^m$.
It carries the spherical class of square $-2$
into a global torus. The same intersection and
asphericity contradiction proves (A.BP.4).
\end{proof}

The condition $P(1,0)\ne0$ is used to produce the
distinct collapsing circles. The exact product
model (A.U.6), for which $P(X,V)=XV(X-V)$, has
$P(1,0)=0$ and is outside this statement.
The variable-coefficient extension uses only the
specific comparison (A.BP.5)--(A.BP.8), with its
unchanged boundary critical curves. The weighted scaling,
simple-zero and endpoint-square identities are
checked in \path{code/E_binary_products.py}.

\paragraph{Two exact scalar-zero comparisons.}
The homogeneous germ
\[
                  h=xy\{z(x-y)+xy\}
\tag{A.BP.9}
\]
has a cubic normal tensor along $(x,y)=0$ which
vanishes identically at $z=0$. On $xy\ne0$ put
\[
 v=x-y,\qquad Z=y(y+v)z+y^2(2y+v).
\]
These are actual coordinates with $y$, with inverse
$x=y+v$, $z=(Z-y^2(2y+v))/(y(y+v))$, and
\[
                              h=vZ+y^4.
\tag{A.BP.10}
\]
Choose an arc between two roots of $y^4=1$ avoiding
$y=0$ and the other roots. The matching circles
$v=\rho A e^{\mathrm i\theta}$,
$Z=\rho^{-1}B e^{-\mathrm i\theta}$, where
$AB=1-y^4$ and $|A|=|B|=|1-y^4|^{1/2}$,
collapse at its endpoints. Choose the positive
constant $\rho$ small enough that $|v|<|y|/2$
on the arc. The entire spherical cycle then lies
in the indicated coordinate chart. Its two endpoint
linkings are $-1$, as in (A.BP.3), so its square
is $-2$. Ordinary homothety in $(x,y,z)$ shrinks
it because (A.BP.9) has degree four. The same
weighted unit map as (A.BP.2), with all three
weights equal to one, includes every local unit
and every positive full-fibre multiplicity.
Thus this particular scalar-zero germ is excluded.

The following different scalar-zero model records
the limit of that argument:
\[
                     h=xz(y^2-x^2z).
\tag{A.BP.11}
\]
For $h=t\ne0$, the substitution $u=y/x$ is
invertible and gives $h=x^3 z(u^2-z)$.
The map forgetting $x$ is a finite etale triple
cover of $E=\{z(u^2-z)\ne0\}$.
There is an explicit finite etale double cover
\[
 \C^*\times(\C\setminus\{1,-1\})\longrightarrow E,
                  \qquad(s,v)\longmapsto(s^2,sv)
\tag{A.BP.12}
\]
in coordinates $(z,u)$, with deck involution
$(s,v)\mapsto(-s,-v)$. Its source is aspherical,
so both $E$ and the regular affine fibre of
(A.BP.11) are aspherical by these actual covering
maps. The quadratic surface branch has normalization
$(v,u)\mapsto(v,vu,u^2)$ and conductor
$2(v=0)+2(u=0)$ in the full reduced surface.
The other two normalization branches have conductors
$2(y=0)+(z=0)$ on $x=0$, and
$(x=0)+2(y=0)$ on $z=0$.
Thus their reduced conductor supports have only
normal-crossing pairs. This is an exact local
model, not a construction of a compact $b_2=0$
fibre. It shows why a vanishing cubic coefficient
alone does not justify a local spherical-cycle
exclusion. The identities and the two finite-cover
formulas are checked in \path{code/E_cubic_double.py}.

